% Options for packages loaded elsewhere
\PassOptionsToPackage{unicode}{hyperref}
\PassOptionsToPackage{hyphens}{url}
\documentclass[
  ukrainian,
]{article}
\usepackage{xcolor}
\usepackage[margin=2.5cm]{geometry}
\usepackage{amsmath,amssymb}
\setcounter{secnumdepth}{-\maxdimen} % remove section numbering
\usepackage{iftex}
\ifPDFTeX
  \usepackage[T1]{fontenc}
  \usepackage[utf8]{inputenc}
  \usepackage{textcomp} % provide euro and other symbols
\else % if luatex or xetex
  \usepackage{unicode-math} % this also loads fontspec
  \defaultfontfeatures{Scale=MatchLowercase}
  \defaultfontfeatures[\rmfamily]{Ligatures=TeX,Scale=1}
\fi
\usepackage{lmodern}
\ifPDFTeX\else
\setmainfont{DejaVu Serif}
\setsansfont{DejaVu Sans}
\setmonofont{DejaVu Sans Mono}
\setmathfont{STIX Two Math}
\fi
\ifPDFTeX\else
  % xetex/luatex font selection
\fi
% Use upquote if available, for straight quotes in verbatim environments
\IfFileExists{upquote.sty}{\usepackage{upquote}}{}
\IfFileExists{microtype.sty}{% use microtype if available
  \usepackage[]{microtype}
  \UseMicrotypeSet[protrusion]{basicmath} % disable protrusion for tt fonts
}{}
\makeatletter
\@ifundefined{KOMAClassName}{% if non-KOMA class
  \IfFileExists{parskip.sty}{%
    \usepackage{parskip}
  }{% else
    \setlength{\parindent}{0pt}
    \setlength{\parskip}{6pt plus 2pt minus 1pt}}
}{% if KOMA class
  \KOMAoptions{parskip=half}}
\makeatother
\usepackage{graphicx}
\graphicspath{{/Users/kalinicenkovitalijmikolajovic/.codex/worktrees/5c15/Script_for_translate_book/Output_Final/}}
\setkeys{Gin}{width=0.74\linewidth,height=0.78\textheight,keepaspectratio}
\makeatletter
\newsavebox\pandoc@box
\newcommand*\pandocbounded[1]{% scales image to fit in text height/width
  \sbox\pandoc@box{#1}%
  \Gscale@div\@tempa{\textheight}{\dimexpr\ht\pandoc@box+\dp\pandoc@box\relax}%
  \Gscale@div\@tempb{\linewidth}{\wd\pandoc@box}%
  \ifdim\@tempb\p@<\@tempa\p@\let\@tempa\@tempb\fi% select the smaller of both
  \ifdim\@tempa\p@<\p@\scalebox{\@tempa}{\usebox\pandoc@box}%
  \else\usebox{\pandoc@box}%
  \fi%
}
% Set default figure placement to htbp
\def\fps@figure{htbp}
\makeatother
\ifLuaTeX
\usepackage[bidi=basic,shorthands=off,provide=*]{babel}
\else
\usepackage[bidi=default,shorthands=off,provide=*]{babel}
\fi
\ifLuaTeX
  \usepackage{selnolig} % disable illegal ligatures
\fi
\setlength{\emergencystretch}{3em} % prevent overfull lines
\providecommand{\tightlist}{%
  \setlength{\itemsep}{0pt}\setlength{\parskip}{0pt}}
\usepackage{bookmark}
\IfFileExists{xurl.sty}{\usepackage{xurl}}{} % add URL line breaks if available
\urlstyle{same}
\hypersetup{
  pdflang={uk},
  hidelinks,
  pdfcreator={LaTeX via pandoc}}

\author{}
\date{}

\begin{document}

\pandocbounded{\includegraphics[keepaspectratio]{images/_page_150_Picture_5.jpeg}}

Фіг. 20.3. Стоячи \(u_n(x, t) = X_n(x)T_n(t)\) .

У нашій задачі рівняння і граничні умови є лінійними та однорідними,
тому довільна сума елементарних стоячих хвиль також є розв'язком
хвильового рівняння і задовольняє граничним умовам. Тепер ми повинні
знайти таку суперпозицію стоячих хвиль, щоб вона також задовольняла
початкові умови, і тоді отримаємо розв'язок нашої проблеми. Заміна
кількості

Зверніть увагу на зображення.

\pandocbounded{\includegraphics[keepaspectratio]{images/_page_154_Figure_5.jpeg}}

Головна різниця між поперечними коливаннями тонкої балки та поперечними
коливаннями струни полягає в тому, що промінь чинить опір згину. Не
заглиблюючись у механіку тонких балок, припустимо, що, враховуючи опір
згинання (замість рівняння хвилі) веде до рівняння четвертого порядку

(21.1)

\section{Шарнірно оперта
балка}\label{ux448ux430ux440ux43dux456ux440ux43dux43e-ux43eux43fux435ux440ux442ux430-ux431ux430ux43bux43aux430}

Розглянемо малі коливання тонкої балки, кінці якої вільно спираються на
дві опори. Це означає, що кінці балки не зміщуються, але можуть вільно
повертатися.

Тому на кінцях балки мають виконуватися граничні умови

\[
u(0, t) = 0, \qquad u(1, t) = 0,
\]

а також умови нульового згинального моменту

\[
u_{xx}(0, t) = 0, \qquad u_{xx}(1, t) = 0.
\]

Із теорії тонких балок випливає, що згинальний момент пропорційний
\(u_{xx}\), тому для вільно опертих кінців саме ці умови є природними.

\pandocbounded{\includegraphics[keepaspectratio]{images/_page_155_Figure_8.jpeg}}

Рис. 21.1. Шарнірно оперта балка.

Отже, коливання балки описуються такою змішаною задачею (для простоти
відповідний коефіцієнт вважаємо рівним одиниці):

(21.2)

\begin{verbatim}
\begin{aligned}
(\text{УЧП})\quad & u_{tt} = -u_{xxxx}, && 0 < x < 1,\ 0 < t < \infty, \\
(\text{ГУ})\quad & u(0,t) = 0, \quad u_{xx}(0,t) = 0, \quad u(1,t) = 0, \quad u_{xx}(1,t) = 0, && 0 < t < \infty, \\
(\text{НУ})\quad & u(x,0) = f(x), \quad u_t(x,0) = g(x), && 0 \leqslant x \leqslant 1.
\end{aligned}
\end{verbatim}

Щоб розв'язати цю задачу, ми використаємо метод розділення змінних. Ми
будемо шукати лише періодичні розв'язки, тобто флуктуації

(21.3)

\[
[2/\pi\omega]^{1/2}
\]

Зверніть увагу, що вибір розв'язку вигляду (21.3) фактично означає, що
ми взяли константу розділення у методі розділення змінних як від'ємну.

Подставим теперь (21.3) в уравнение колебаний балки и получим уравнение
для \(X\left(x\right)\)

\[
\frac{2}{\pi\omega} [1-\cos(\omega a)]
\]

общее решение которого записывается в виде

1

Щоб визначити коефіцієнти C, D, E та F, підставляємо загальний розв'язок
у граничні умови

\[
e^{-1}a + \omega
\frac{1}{2}e^{-\alpha}\sin\omega
\begin{split} u\left(0,\ t\right) &= 0 \Rightarrow X\left(0\right)T\left(t\right) = 0 \Rightarrow X\left(0\right) = 0 \Rightarrow C + E = 0 \\ u_{xx}\left(0,\ t\right) &= 0 \Rightarrow X''\left(0\right)T\left(t\right) = 0 \Rightarrow X''\left(0\right) = 0 \Rightarrow -C + E = 0 \\ u\left(1,\ t\right) &= 0 \Rightarrow D\sin\sqrt{\omega} + F\sin\sqrt{\omega} = 0, \\ u_{xx}\left(1,\ t\right) &= 0 \Rightarrow -D\sin\sqrt{\omega} + F\sin\sqrt{\omega} = 0. \end{split}
\frac{1-e^{-a\omega}}{\omega}
\frac{2}{\pi}F''(\omega)
F \sinh \sqrt{\omega} = 0,
f(x) = \int_{0}^{\infty} F(\omega) \cos(\omega x) d\omega \qquad \qquad F(x) = \frac{2}{\pi} \int_{0}^{\infty} f(x) \cos(\omega x) dx
D \sin \sqrt{\omega} = 0.
\]

\begin{enumerate}
\def\labelenumi{\arabic{enumi}.}
\tightlist
\item
\end{enumerate}

\[
F = 0
\]

\section{ПЕРЕХІД ДО БЕЗРОЗМІРНИХ
ЗМІННИХ}\label{ux43fux435ux440ux435ux445ux456ux434-ux434ux43e-ux431ux435ux437ux440ux43eux437ux43cux456ux440ux43dux438ux445-ux437ux43cux456ux43dux43dux438ux445}

МЕТА ЛЕКЦІЇ: Показати, яким чином крайові, початкові та інші фізичні
задачі можна записувати у безрозмірному вигляді. У різних галузях знань
рівняння формулюються по-різному, але після переходу до безрозмірних
змінних їх часто можна звести до спільної математичної форми.

Переход к безразмерным переменным позволяет привести уравнения к одному
и тому же виду. Именно по этой причине при изучении уравнений с частными
производными в математике стараются отвлекаться от физического смысла
входящих в них параметров. Вот почему химик, физик или биолог должен
преобразовать свое уравнение к той форме, которая принята в нашей книге.

Основная идея введения новых (безразмерных) переменных состоит в том,
что после перехода к безразмерным переменным вадача становится чисто
математической и уже не содержит

карактерных физических констант. Именно таким способом многие различные
уравнения физики, биологии и химии, содержащие всякие нюансы, связанные
с физическими параметрами, приводятся к одной и той же простой форме
(см. рис. 22.1).

Познакомимся с этим преобразованием на простом примере.

\pandocbounded{\includegraphics[keepaspectratio]{images/_page_161_Figure_4.jpeg}}

Рис. 22.1. Некоторые задачи, приводящиеся к одной и той же безразмерной
форме.

\section{Приведення дифузійної задачі до безрозмірного
вигляду}\label{ux43fux440ux438ux432ux435ux434ux435ux43dux43dux44f-ux434ux438ux444ux443ux437ux456ux439ux43dux43eux457-ux437ux430ux434ux430ux447ux456-ux434ux43e-ux431ux435ux437ux440ux43eux437ux43cux456ux440ux43dux43eux433ux43e-ux432ux438ux433ux43bux44fux434ux443}

\begin{verbatim}
Начнем со смешанной задачи для отрезка с начальным условием на температуру  u\left(x,\,0\right)=\sin\left(\pi x/L\right)  и граничными температурами на концах  T_1  и  T_2 . Другими словами, мы решаем задачу
\end{verbatim}

(22.1)

\begin{verbatim}
\begin{array}{ccccccccccccccccccccccccccccccccccc
\end{verbatim}

(см. рис. 22.2).

\pandocbounded{\includegraphics[keepaspectratio]{images/_page_161_Figure_10.jpeg}}

Рис. 22.2. Задача теплопроводности в плоскости переменных (х, т).

Наша цель --- дать новую эквивалентную (22.1) постановку задачи таким
образом, чтобы

в новой задаче не было никаких физических параметров (вроде α);

\begin{enumerate}
\def\labelenumi{\arabic{enumi})}
\setcounter{enumi}{1}
\tightlist
\item
  начальные и граничные условия стали проще.
\end{enumerate}

\begin{verbatim}
Для осуществления этой цели, мы введем три новых безразмерных переменных U,   \xi   и   \tau  , чтобы заменить прежние u, x и t по схеме
\end{verbatim}

\[
u \to U
\]

\begin{verbatim}
(безразмерная температура),   x \to \xi   (безразмерная длина),   t \to \tau   (безразмерное время).
\end{verbatim}

Мы проведем все три преобразования одновременно.

\begin{verbatim}
# Перетворення залежної змінної  u \rightarrow U
\end{verbatim}

Определим функцию \(U(x, t)\) по формуле

\[
U(x, t) = \frac{u(x, t) - T_1}{T_2 - T_1}
\]

.

\begin{verbatim}
Ясно, что новая температура   U\left(x,\,t\right)   безразмерна, поскольку мы делим градусы Цельсия на градусы Цельсия. Новые граничные условия для   U\left(x,\,t\right)   при x=0 и x=L примут вид   U\left(0,\,t\right)=0   и   U\left(L,\,t\right)=1  . Давайте поставим задачу для новой функции   U\left(x,\,t\right)  . После несложных преобразований задача (22.1) переходит в новую задачу
\end{verbatim}

задачу

\begin{verbatim}
\begin{array}{ll} \text{(УЧП)} & U_t = \alpha^2 U_{xx}, \quad 0 < x < L, \quad 0 < t < \infty, \\ U_t = 0, \quad 0 < t < \infty, \\ U_t = 0, \quad 0 < t < \infty, \\ U_t = 0, \quad 0 < t < \infty, \\ U_t = 0, \quad 0 < t < \infty, \\ U_t = 0, \quad 0 < t < \infty, \\ U_t = 0, \quad 0 < t < \infty, \\ U_t = 0, \quad 0 < t < \infty, \\ U_t = 0, \quad 0 < t < \infty, \\ U_t = 0, \quad 0 < t < \infty, \\ U_t = 0, \quad 0 < t < \infty, \\ U_t = 0, \quad 0 < t < \infty, \\ U_t = 0, \quad 0 < t < \infty, \\ U_t = 0, \quad 0 < t < \infty, \\ U_t = 0, \quad 0 < t < \infty, \\ U_t = 0, \quad 0 < t < \infty, \\ U_t = 0, \quad 0 < t < \infty, \\ U_t = 0, \quad 0 < t < \infty, \\ U_t = 0, \quad 0 < t < \infty, \\ U_t = 0, \quad 0 < t < \infty, \\ U_t = 0, \quad 0 < t < \infty, \\ U_t = 0, \quad 0 < t < \infty, \\ U_t = 0, \quad 0 < t < \infty, \\ U_t = 0, \quad 0 < t < \infty, \\ U_t = 0, \quad 0 < t < \infty, \\ U_t = 0, \quad 0 < t < \infty, \\ U_t = 0, \quad 0 < t < \infty, \\ U_t = 0, \quad 0 < t < \infty, \\ U_t = 0, \quad 0 < t < \infty, \\ U_t = 0, \quad 0 < t < \infty, \\ U_t = 0, \quad 0 < t < \infty, \\ U_t = 0, \quad 0 < t < \infty, \\ U_t = 0, \quad 0 < t < \infty, \\ U_t = 0, \quad 0 < t < \infty, \\ U_t = 0, \quad 0 < t < \infty, \\ U_t = 0, \quad 0 < t < \infty, \\ U_t = 0, \quad 0 < t < \infty, \\ U_t = 0, \quad 0 < t < \infty, \\ U_t = 0, \quad 0 < t < \infty, \\ U_t = 0, \quad 0 < t < \infty, \\ U_t = 0, \quad 0 < t < \infty, \\ U_t = 0, \quad 0 < t < \infty, \\ U_t = 0, \quad 0 < t < \infty, \\ U_t = 0, \quad 0 < t < \infty, \\ U_t = 0, \quad 0 < t < \infty, \\ U_t = 0, \quad 0 < t < \infty, \\ U_t = 0, \quad 0 < t < \infty, \\ U_t = 0, \quad 0 < t < \infty, \\ U_t = 0, \quad 0 < t < \infty, \\ U_t = 0, \quad 0 < t < \infty, \\ U_t = 0, \quad 0 < t < \infty, \\ U_t = 0, \quad 0 < t < \infty, \\ U_t = 0, \quad 0 < t < \infty, \\ U_t = 0, \quad 0 < t < \infty, \\ U_t = 0, \quad 0 < t < \infty, \\ U_t = 0, \quad 0 < t < \infty, \\ U_t = 0, \quad 0 < t < \infty, \\ U_t = 0, \quad 0 < t < \infty, \\ U_t = 0, \quad 0 < t < \infty, \\ U_t = 0, \quad 0 < t < \infty, \\ U_t = 0, \quad 0 < t < \infty, \\ U_t = 0, \quad 0 < t < \infty, \\ U_t = 0, \quad 0 < t < \infty, \\ U_t = 0, \quad 0 < t < \infty, \\ U_t = 0, \quad 0 < t < \infty, \\ U_t = 0, \quad 0 < t < \infty, \\ U_t = 0, \quad 0 < t < \infty, \\ U_t = 0, \quad 0 < t < \infty, \\ U_t = 0, \quad 0 < t < \infty, \\ U_t = 0, \quad 0 < t < \infty, \\ U_t = 0, \quad 0 < t < \infty, \\ U_t = 0, \quad 0 < t < \infty, \\ U_t = 0, \quad 0 < t < \infty, \\ U_t = 0, \quad 0 < t < \infty, \\ U_t = 0, \quad 0 < t < \infty, \\ U_t = 0, \quad 0 < t < \infty, \\ U_t = 0, \quad 0 < t < \infty, \\ U_t = 0, \quad 0 < t < \infty, \\ U_t = 0, \quad 0 < t < \infty, \\ U_t = 0, \quad 0 < t < \infty, \\ U_t = 0, \quad 0 < t < \infty, \\ U_t = 0, \quad 0 < t < \infty, \\ U_t = 0, \quad 0 < t < \infty, \\ U_t
\end{verbatim}

\begin{verbatim}
Если угодно, то здесь можно остановиться, решить задачу для   U\left(x,\, t\right)   и затем найти   u\left(x,\,t\right)   по формуле
\end{verbatim}

\[
u(x, t) = T_1 + (T_2 - T_1) U(x, t).
\]

Однако мы пойдем дальше и займемся преобразованием независимых
переменных x и t.

\begin{verbatim}
# Перетворення просторової координати  x \rightarrow \xi
\end{verbatim}

\begin{verbatim}
Представляется совершенно очевидным, как следует выбрать безразмерную переменную   \xi  . Поскольку   0 \leqslant x \leqslant L  , полагаем
\end{verbatim}

\[
\xi = x/L
\]

.

После вычисления производных

\[
U_x = U_{\xi} \xi_x = \frac{1}{L} U_{\xi},
U_{xx} = \frac{1}{L^2} U_{\xi\xi}
\]

наступна задача стає очевидною (у змінних \(U\), \(\xi\), \(t\)):

\begin{verbatim}
\begin{aligned}
(\text{УЧП})\quad & U_t = (\alpha/L)^2 U_{\xi\xi}, && 0 < \xi < 1,\quad 0 < t < \infty, \\
(\text{ГУ})\quad & U(0,t)=0,\quad U(1,t)=1, && 0 < t < \infty, \\
(\text{НУ})\quad & U(\xi,0)=\frac{\sin(\pi\xi)-T_1}{T_2-T_1}, && 0 \leqslant \xi \leqslant 1.
\end{aligned}
\end{verbatim}

Ми вже пройшли дві третини шляху. Останній крок пов'язаний із введенням
безрозмірного часу \(\tau\) для виключення коефіцієнта \([\alpha/L]^2\)
з диференціального рівняння.

\section{\texorpdfstring{\(t \rightarrow \tau\) Перетворення
часу}{t \textbackslash rightarrow \textbackslash tau Перетворення часу}}\label{t-rightarrow-tau-ux43fux435ux440ux435ux442ux432ux43eux440ux435ux43dux43dux44f-ux447ux430ux441ux443}

Як обрати час без розмірів? Це не так очевидно, як у випадку перших двох
змінних. Однак, оскільки наша мета --- виключити сталу \([\alpha/L]^2\)
з рівняння (22.3), ми діятимемо наступним чином.

\begin{itemize}
\item
  \begin{enumerate}
  \def\labelenumi{\arabic{enumi}.}
  \tightlist
  \item
    Спробуємо замінити вигляд \(\tau = ct\) , де c --- невідома
    константа.
  \end{enumerate}

  \begin{itemize}
  \tightlist
  \item
    \begin{enumerate}
    \def\labelenumi{\arabic{enumi}.}
    \setcounter{enumi}{1}
    \tightlist
    \item
      Давайте розрахуємо \(u_t = u_{\tau} \tau_t = c u_{\tau}\) .
    \end{enumerate}
  \item
    \begin{enumerate}
    \def\labelenumi{\arabic{enumi}.}
    \setcounter{enumi}{2}
    \tightlist
    \item
      Підставимо цю похідну у рівняння (22.3) і отримаємо
    \end{enumerate}
  \end{itemize}
\end{itemize}

\[
cu_{\mathbf{t}} = [\alpha/L]^2 u_{\xi\xi}.
\]

Тому потрібно обрати \(c = [\alpha/L]^2\) , тобто новий час слід вводити
за формулою

\[
\tau = [\alpha/L]^2 t.
\]

Застосовуючи це перетворення до попередньої задачі (22.3), ми повністю
переходимо до безрозмірних змінних \(U\), \(\xi\) та \(\tau\):

\begin{verbatim}
\begin{aligned}
(\text{УЧП})\quad & U_{\tau} = U_{\xi\xi}, && 0 < \xi < 1,\quad 0 < \tau < \infty, \\
(\text{ГУ})\quad & U(0,\tau) = 0,\quad U(1,\tau) = 1, && 0 < \tau < \infty, \\
(\text{НУ})\quad & U(\xi,0) = \frac{\sin(\pi \xi) - T_1}{T_2 - T_1}, && 0 \le \xi \le 1.
\end{aligned}
\end{verbatim}

Нова безрозмірна система має такі властивості:

\begin{enumerate}
\def\labelenumi{\arabic{enumi}.}
\tightlist
\item
  Рівняння не містить фізичних параметрів.
\item
  Граничні умови мають простий вигляд.
\item
  Початкова умова суттєво не ускладнюється і містить лише відомі
  функції.
\item
  Якщо знайдено розв'язок \(U(\xi,\tau)\), то розв'язок початкової
  задачі відновлюється за формулою
\end{enumerate}

\[
u(x,t) = T_1 + (T_2 - T_1)\,U\left(x/L,\alpha^2 t/L^2\right).
\]

На цьому завершується обговорення безрозмірного запису задач.
Універсального рецепта введення нових змінних не існує: зазвичай
доводиться спиратися на фізичну інтуїцію і перевіряти кілька природних
варіантів.

Завершимо лекцію простим прикладом.

\section{Приклад перетворення гіперболічної задачі на безрозмірний
вигляд}\label{ux43fux440ux438ux43aux43bux430ux434-ux43fux435ux440ux435ux442ux432ux43eux440ux435ux43dux43dux44f-ux433ux456ux43fux435ux440ux431ux43eux43bux456ux447ux43dux43eux457-ux437ux430ux434ux430ux447ux456-ux43dux430-ux431ux435ux437ux440ux43eux437ux43cux456ux440ux43dux438ux439-ux432ux438ux433ux43bux44fux434}

Розглянемо коливну струну.

(22.5)

\begin{verbatim}
\begin{aligned}
(\text{УЧП})\quad & u_{tt} = \alpha^2 u_{xx}, && 0 < x < L,\quad 0 < t < \infty, \\
(\text{ГУ})\quad & u(0,t) = 0,\quad u(L,t) = 0, && 0 < t < \infty, \\
(\text{НУ})\quad & u(x,0) = \sin(\pi x/L) + 0.5\sin(3\pi x/L),\quad u_t(x,0) = 0, && 0 \le x \le L.
\end{aligned}
\end{verbatim}

Переходячи до безрозмірних змінних

\[
\xi = x/L, \qquad \tau = (\alpha/L)t,
\]

одержуємо нову задачу

\begin{verbatim}
\begin{aligned}
(\text{УЧП})\quad & u_{\tau\tau} = u_{\xi\xi}, && 0 < \xi < 1,\quad 0 < \tau < \infty, \\
(\text{ГУ})\quad & u(0,\tau) = 0,\quad u(1,\tau) = 0, && 0 < \tau < \infty, \\
(\text{НУ})\quad & u(\xi,0) = \sin(\pi \xi) + 0.5\sin(3\pi \xi),\quad u_{\tau}(\xi,0) = 0, && 0 \le \xi \le 1.
\end{aligned}
\end{verbatim}

Розв'язком цієї задачі є функція

\[
u(\xi,\tau) = \cos(\pi \tau)\sin(\pi \xi) + 0.5\cos(3\pi \tau)\sin(3\pi \xi).
\]

Повертаючись до координат \(x\) і \(t\), дістаємо

\[
u(x,t) = \cos(\pi \alpha t/L)\sin(\pi x/L) + 0.5\cos(3\pi \alpha t/L)\sin(3\pi x/L).
\]

\subsection{Зауваження}\label{ux437ux430ux443ux432ux430ux436ux435ux43dux43dux44f}

\begin{enumerate}
\def\labelenumi{\arabic{enumi}.}
\tightlist
\item
  Аналіз розмірностей особливо важливий у чисельних методах, оскільки
  більшість програм орієнтована на загальні математичні моделі, а не на
  конкретний набір фізичних параметрів.
\item
  Не завжди потрібно робити безрозмірними всі змінні; інколи достатньо
  перетворити лише одну або дві з них.
\item
  Безрозмірний запис часто дозволяє побачити головні параметри задачі
  значно чіткіше.
\end{enumerate}

\section{ЗАВДАННЯ}\label{ux437ux430ux432ux434ux430ux43dux43dux44f}

\begin{enumerate}
\def\labelenumi{\arabic{enumi}.}
\tightlist
\item
  За допомогою заміни
\end{enumerate}

\[
\xi = x/L, \qquad \tau = (\alpha/L)t
\]

перетворіть задачу (22.5) до безрозмірного вигляду.

\begin{enumerate}
\def\labelenumi{\arabic{enumi}.}
\setcounter{enumi}{1}
\tightlist
\item
  Знайдіть безрозмірну форму задачі
\end{enumerate}

\begin{verbatim}
\begin{aligned}
(\text{УЧП})\quad & u_t = \alpha^2 u_{xx}, && 0 < x < L, \\
(\text{ГУ})\quad & u(0,t) = T_1,\quad u(L,t) = 0, && 0 < t < \infty, \\
(\text{НУ})\quad & u(x,0) = T_2, && 0 \le x \le L.
\end{aligned}
\end{verbatim}

\begin{enumerate}
\def\labelenumi{\arabic{enumi}.}
\setcounter{enumi}{2}
\tightlist
\item
  За допомогою заміни
\end{enumerate}

\[
U(x,t) = \frac{u(x,t) - T_1}{T_2 - T_1}
\]

перетворіть задачу (22.1) на задачу (22.2).

\begin{enumerate}
\def\labelenumi{\arabic{enumi}.}
\setcounter{enumi}{3}
\item
  Поясніть, чому заміна змінної \(\tau = \alpha t\) усуває параметр
  \(\alpha\) з хвильового рівняння \(u_{tt} = \alpha^2 u_{xx}\) і який
  фізичний зміст має ця заміна.
\item
  Як слід вибрати нову просторову координату \(\xi\), щоб усунути
  параметр \(v\) з рівняння
\end{enumerate}

\[
u_t + vu_x = 0?
\]

\subsection{Лекція 23}\label{ux43bux435ux43aux446ux456ux44f-23}

\section{КЛАСИФІКАЦІЯ РІВНЯНЬ У ЧАСТИННИХ ПОХІДНИХ (КАНОНІЧНА ФОРМА
ГІПЕРБОЛІЧНОГО
РІВНЯННЯ)}\label{ux43aux43bux430ux441ux438ux444ux456ux43aux430ux446ux456ux44f-ux440ux456ux432ux43dux44fux43dux44c-ux443-ux447ux430ux441ux442ux438ux43dux43dux438ux445-ux43fux43eux445ux456ux434ux43dux438ux445-ux43aux430ux43dux43eux43dux456ux447ux43dux430-ux444ux43eux440ux43cux430-ux433ux456ux43fux435ux440ux431ux43eux43bux456ux447ux43dux43eux433ux43e-ux440ux456ux432ux43dux44fux43dux43dux44f}

МЕТА ЛЕКЦІЇ: Показати, що лінійне рівняння другого порядку з двома
незалежними змінними

\[
Au_{xx} + Bu_{xy} + Cu_{yy} + Du_x + Eu_y + Fu = G
\]

(A, B, C, D, E, F и G являются функциями x и y или могут быть
константами) относится к одному из следующих типов:

\begin{itemize}
\tightlist
\item
  \begin{enumerate}
  \def\labelenumi{\arabic{enumi})}
  \tightlist
  \item
    гиперболическому (если \(B^2 4AC > 0\) ),
  \end{enumerate}
\item
  \begin{enumerate}
  \def\labelenumi{\arabic{enumi})}
  \setcounter{enumi}{1}
  \tightlist
  \item
    параболическому (если \(B^2 4AC =\) ),
  \end{enumerate}
\item
  \begin{enumerate}
  \def\labelenumi{\arabic{enumi})}
  \setcounter{enumi}{2}
  \tightlist
  \item
    эллиптическому (если \(0 B^2 4AC < 0\) ).
  \end{enumerate}
\end{itemize}

Показать, какие новые переменные \(\xi = \xi(x, y)\) и
\(\eta = \eta(x, y)\) можно ввести вместо x и y, чтобы упростить
исходное уравнение. В новых переменных \(\xi\) и \(\eta\) уравнение с
частными производными приводится к одному из следующих видов (в
вависимости от того, будет ли величина \(B^2 - 4AC\) положительна, равна
нулю или отрицательна):

\begin{itemize}
\tightlist
\item
  \begin{enumerate}
  \def\labelenumi{\arabic{enumi}.}
  \tightlist
  \item
  \end{enumerate}
\end{itemize}

\[
\begin{cases} u_{\xi\xi} u_{\eta\eta} = \Psi(\xi,  \eta, u, u_{\xi}, u_{\eta}) і \text{(дві канонічні форми гіперболічного рівняння),} \end{cases}
\]

\begin{itemize}
\tightlist
\item
  \begin{enumerate}
  \def\labelenumi{\arabic{enumi}.}
  \setcounter{enumi}{1}
  \tightlist
  \item
    \(u_{\eta\eta} = \Phi(\xi, \eta, u, u_{\xi}, u_{\eta})\)
    (каноническая форма параболического уравнения),
  \end{enumerate}
\item
  \begin{enumerate}
  \def\labelenumi{\arabic{enumi}.}
  \setcounter{enumi}{2}
  \tightlist
  \item
    \(u_{\xi\xi} + u_{\eta\eta} = \Phi(\xi, \eta, u, u_{\xi}, u_{\eta})\)
    (каноническая форма эллиптического уравнения),
  \end{enumerate}
\end{itemize}

где \(\Phi\) и \(\Psi\) ---функции от частных производных первого
порядка, зависимой переменной u и новых независимых переменных \(\xi\) и
\(\eta\) . Конкретный вид функций \(\Phi\) и \(\Psi\) зависит от
исходного уравнения.

У читателя может возникнуть мысль, что глава, посвященная классификации
уравнений, должна размещаться в самом начале книги. Вероятно, это
правильно, и многие книги действительно начинаются с этого материала.
Однако правда и то, что многие студенты не могут с энтузиазмом изучать
нечто такое, о чем у них нет никакого представления. Именно по этой
причине мы отложили изучение темы, посвященной классификации уравнений с
частными производными, до настоящего момента.

Наша задача - классифицировать все уравнения вида

(23 1)

\[
Au_{xx} + Bu_{xy} + Cu_{yy} + Du_x + Eu_y + Fu = G,
\]

де \(A\), \(B\), \(C\), \(D\), \(E\), \(F\) і \(G\) є, взагалі кажучи,
функціями \(x\) та \(y\). Згідно з цією схемою рівняння називається

\begin{enumerate}
\def\labelenumi{\arabic{enumi}.}
\tightlist
\item
  гіперболічним у точці \((x_0, y_0)\), якщо
  \(B^2(x_0, y_0) - 4A(x_0, y_0)C(x_0, y_0) > 0\);
\item
  параболічним у точці \((x_0, y_0)\), якщо
  \(B^2(x_0, y_0) - 4A(x_0, y_0)C(x_0, y_0) = 0\);
\item
  еліптичним у точці \((x_0, y_0)\), якщо
  \(B^2(x_0, y_0) - 4A(x_0, y_0)C(x_0, y_0) < 0\).
\end{enumerate}

Залежно від типу його можна звести до відповідної канонічної форми. Щоб
краще зрозуміти класифікаційну схему, наведемо кілька прикладів.

\section{Приклади гіперболічних, параболічних та еліптичних
рівнянь}\label{ux43fux440ux438ux43aux43bux430ux434ux438-ux433ux456ux43fux435ux440ux431ux43eux43bux456ux447ux43dux438ux445-ux43fux430ux440ux430ux431ux43eux43bux456ux447ux43dux438ux445-ux442ux430-ux435ux43bux456ux43fux442ux438ux447ux43dux438ux445-ux440ux456ux432ux43dux44fux43dux44c}

\begin{enumerate}
\def\labelenumi{\arabic{enumi}.}
\tightlist
\item
  У равнение теплопроводности \(u_t = u_{xx}\) --- линейное уравнение
  второго порядка вида (23.1) с коэффициентами
\end{enumerate}

\[
A=1
\]

, \(B=0\) , \(C=0\) , \(D=0\) , \(E=-1\) , \(F=0\) , \(G=0\) ,

так что \(B^2-4AC=0\) при всех x и t. Следовательно, это уравнение будет
параболическим при всех x и t. В общем уравнении (23 1) переменная y
играет роль времени t. Результат не изменится, если в общем уравнении
поменять местами x и y.

\begin{enumerate}
\def\labelenumi{\arabic{enumi}.}
\setcounter{enumi}{1}
\tightlist
\item
  Волновое уравнение \(u_{1t} = u_{xx}\) также принадлежит классу
  уравнений вида (23.1) с коэффициентами
\end{enumerate}

\[
A=1
\]

, \(B=0\) , \(C=-1\) , \(D=E=F=G=0\) .

Так как \(B^2 - 4AC = 4\) для всех x и t, то это уравнение является
еиперболическим при всех x и t.

1) Следует указать, что приведенные в примерах уравнения являются
частными случаями более общих параболических, гиперболических в
эллиптических уравнений, а основные их свойства характерны и для общего
случая.

\begin{enumerate}
\def\labelenumi{\arabic{enumi}.}
\setcounter{enumi}{2}
\item
  У равнение Лапласа \(u_{xx} + u_{yy} = 0\) является эллиптическим при
  всех x и y, так как \(B^2 - 4AC = - 4 < 0\) .
\item
  Линейное уравнение \(xu_{xx} + u_{yy} = \sin x\) с переменными
  коэффициентами также принадлежит классу уравнений вида (23.1), но для
  него \(B^2-4AC=-4x\) и, следовательно, уравнение будет эллиптическим
  при x \textgreater{} 0,
\end{enumerate}

лараболическим при x=0, гиперболическим при x \textless{} 0.

Этот пример иллюстрирует тот факт, что тип уравнения о переменными
коэффициентами может изменяться от точки к точке.

Читатель должен обратить внимание на следующее: тип уравнения (23.1)
определяется только коэффициентами при вторых првизводных и никак не
зависит ни от коэффициентов при первых производных и самой функции, ни
от свободного члена.

Обратимся теперь к главной теме нашей лекции --- приведению
гиперболических уравнений к каноническому виду. Известно, что если в
данной области пространства уравнение является гиперболическим, то
вместо x и y можно ввести координаты \(\xi\) и \(\eta\)
(характеристические переменные) таким образом, что уравнение примет
простейшую форму

(23.2)

\[
u_{\xi\eta} = \Phi(\xi, \eta, u, u_{\xi}, u_{\eta}).
\]

Это уравнение содержит только одну производную второго порядка
\(u_{\xi\eta}\) , а функция \(\Phi(\xi,) \eta,u,u_{\xi},u_{\eta})\)
зависит только от новых независимых переменных § и η, зависимой
переменной и и первых производных из и ип. Конкретный вид функции Ф
зависит, конечно, от исходного уравнения и формул перехода к новым
координатам § и п. Научимся отыскивать эти функции.

\section{Каноническая форма гиперболического
уравнения}\label{ux43aux430ux43dux43eux43dux438ux447ux435ux441ux43aux430ux44f-ux444ux43eux440ux43cux430-ux433ux438ux43fux435ux440ux431ux43eux43bux438ux447ux435ux441ux43aux43eux433ux43e-ux443ux440ux430ux432ux43dux435ux43dux438ux44f}

Рассмотрим уравнение общего вида

(23.3)

\[
Au_{xx} + Bu_{xy} + Cu_{yy} + Du_x + Eu_y + Fu = G
\]

и пусть в интересующей нас области \(B^2 - 4AC > 0\) . Нам нужно ввести
новые координаты 1)

\[
\xi = \xi(x, y),
\eta = \eta(x, y)
\]

\begin{verbatim}
тому в рівнянні (23.3) залишається лише одна часткова похідна другого порядку  u_{\xi\eta}  (не слід шукати заміну змінним, за допомогою якої гіперболічне вирівнювання
\end{verbatim}

1) Необхідно, щоб це перетворення координат було локально оборотним і
подвійно диференційованим. --- Ед. Ед.

У цьому випадку це неможливо).

Вычислим сначала частные производные

\[
\begin{array}{c} u_x = u_\xi \xi_x + u_\eta \eta_x, \\ u_y = u_\xi \xi_y + u_\eta \eta_y, \\ (23.4) \ u_{xx} = u_{\xi\xi} \xi_x^2 + 2u_{\xi\eta} \xi_x \eta_x + u_{\eta\eta} \eta_x^2 + u_{\xi} \xi_{xx} + u_{\eta} \eta_{xx}, \\ u_{xy} = u_{\xi\xi} \xi_x^2 + u_{\xi\eta} (\xi_x \eta_y + \xi_y \eta_x) + u_{\eta\eta} \eta_x \eta_y + u_{\xi} \xi_{xy} + u_{\eta} \eta_{xy}, \\ u_{yy} = u_{\xi\xi} \xi_y^2 + 2u_{\xi\eta} \xi_y \eta_y + u_{\eta\eta} \eta_y^2 + u_{\xi} \xi_{yy} + u_{\eta} \eta_{yy}. \end{array}
\]

Підставляючи ці відношення до початкового рівняння (23.3), після
простих, але громіздких обчислень отримуємо

(23.5)

\[
\overline{A}u_{\xi\xi} + \overline{B}u_{\xi\eta} + \overline{C}u_{\eta\eta} + \overline{D}u_{\xi} + \overline{E}u_{\eta} + \overline{F}u = \overline{G},
\overline{A} = A\xi_{x}^{2} + B\xi_{x}\xi_{y} + C\xi_{y}^{2},
\overline{B} = 2A\xi_{x}\eta_{y} + B(\xi_{x}\eta_{y} + \xi_{y}\eta_{x}) + 2C\xi_{y}\eta_{y},
\overline{C} = A\eta_{x}^{2} + B\eta_{x}\eta_{y} + C\eta_{y}^{2},
\overline{D} = A\xi_{xx} + B\xi_{xy} + C\xi_{yy} + D\xi_{x} + E\xi_{y},
\overline{E} = A\eta_{xx} + B\eta_{xy} + C\eta_{yy} + D\eta_{x} + E\eta_{y},
\overline{F} = F,
\overline{G} = G.
\]

Читателю предоставится возможность проделать эти выкладки при решении
задач.

Наш следующий шаг---выбрать функции \(\xi = \xi(x, y)\) и
\(\eta = \eta(x, y)\) таким образом, чтобы коэффициенты \(\overline{A}\)
и \(\overline{C}\) обратились в нуль. Это позволит нам привести исходное
уравиение к канонической форме. Итак, необходимо потребовать, чтобы

\[
\bar{A} = A\xi_x^2 + B\xi_x\xi_y + C\xi_y^2 = 0,
\]

\(\bar{C} = A\eta_x^2 + B\eta_x\eta_y + C\eta_y^2 = 0.\) Эти уравнения
можно представить в виде

\[
\begin{split} &A\left[\xi_{x}/\xi_{y}\right]^{2}+B\left[\xi_{x}/\xi_{y}\right]+C=0,\\ &A\left[\eta_{x}/\eta_{y}\right]^{2}+B\left[\eta_{x}/\eta_{y}\right]^{2}+C=0. \end{split}
\]

Разрешив их относительно \([\xi_x/\xi_y]\) и \([\eta_x/\eta_y]\) ,
получаем

(23.7)

\[
\begin{bmatrix} \xi_x/\xi_y \end{bmatrix}
= \frac{-B + \sqrt{B^2 - 4AC}}{2A}
\]

(характеристические

\[
[\eta_x/\eta_y] = \frac{-B - \sqrt{B^2 - 4AC}}{2A}
\]

уравнения).

\begin{verbatim}
Кожне квадратичне рівняння для  [\xi_x/\xi_y]  і  [\eta_x/\eta_y]  має два ядра, але залишаємо лише одне, щоб вони відрізнялися.
\end{verbatim}

\pandocbounded{\includegraphics[keepaspectratio]{images/_page_170_Figure_3.jpeg}}

\begin{verbatim}
Фіг. 23.1. Характеристики  \xi(x, y) = c  та  \eta(x, y) = c  .
\end{verbatim}

\begin{verbatim}
Задача була зведена до знаходження двох функцій  \xi(x, y)  і  \eta(x, y)  , так що відношення  [\xi_x/\xi_y]  і  [\eta_x/\eta_y]  задовольняють рівняння (23.7). Дуже легко знайти такі функції 1) якщо звернути увагу на рис. 23.1. Щоб зрозуміти, як це можна зробити, розглянемо просте рівняння
\end{verbatim}

\[
u_{xx} - 4u_{yy} + u_x = 0.
\]

Его характеристики определяются из уравнений

\begin{verbatim}
\begin{aligned}
\frac{dy}{dx} &= -\frac{\xi_x}{\xi_y} = \frac{B - \sqrt{B^2 - 4AC}}{2A} = -2, \\
\frac{dy}{dx} &= -\frac{\eta_x}{\eta_y} = \frac{B + \sqrt{B^2 - 4AC}}{2A} = 2.
\end{aligned}
\end{verbatim}

Решая эти уравнения относительно у, получаем

\[
y = -2x + c_1,
y = 2x + c_2.
\]

Чтобы найти \(\xi\) и \(\eta\) , разрешим полученные уравнения
относительно констант \(c_1\) и \(c_2\) , оставляя их в правой части и
перенося остальные слагаемые в левую.

Выражения, стоящие в левых частях, примем за искомые функции § и пр

\[
\xi = y + 2x = c_1,
\eta = y - 2x = c_2.
\]

1) См., например, Тихонов А. Н., Самарский А. А. Уравнения
математической физики. --- М.: Наука, 1972 и др. издания. --- Прим. ред.

\begin{verbatim}
Очевидно, що функції  \xi  і  \eta  введені таким чином, задовольняють рівняння характеристик. Ці нові координати показані на рис. 23.2. На цьому наше обговорення того, як знайти нові координати, завершується. Останнє, що залишилося — визначити новий тип початкового рівняння.
\end{verbatim}

\pandocbounded{\includegraphics[keepaspectratio]{images/_page_171_Figure_3.jpeg}}

Фіг. 23.2. Характеристичні координати рівняння \(u_{xx}\) - 4u\_\{uu\} +
\(u_x = 0\).

\begin{verbatim}
Це досить просто: знайти канонічну форму диференціального рівняння з частковими похідними, підставити нові координати  \xi(x, y)  і  \eta(x,) y)  у рівняння
\end{verbatim}

\[
\overline{A}u_{\xi\xi} + \overline{B}u_{\xi\eta} + \overline{C}u_{\eta\eta} + \overline{D}u_{\xi} + \overline{E}u_{\eta} + \overline{F}u = \overline{G}
\]

,

\begin{verbatim}
де  \overline{A}  ,  \overline{B}  ,  \overline{C}  ,  \overline{D}  ,  \overline{E}  ,  \overline{F}  та  \overline{G}  визначаються формулами (23.6).
\end{verbatim}

Прежде чем закончить лекцию, давайте на конкретном примере посмотрим,
как «работает» общий метод приведения к каноническому виду.

Приведение гиперболического уравнения
y\textsuperscript{2u\_\{xx\}-x}2u\_\{yy\}=0 к канонической форме

Рассмотрим уравнение

\[
y^2u_{xx}-x^2u_{yy}=0
,  x>  ,  0 y>0  ,
\]

що є гіперболічним у першому квадранті. Знайдемо нові координати так,
щоб початкове рівняння набуло канонічної форми.

ШАГ 1. Решаем уравнения характеристик

\[
\frac{dy}{dx} = \frac{B + \sqrt{B^2 - 4AC}}{2A} = \frac{x}{y}
\]

(напомним, что этот шаг

\[
\frac{dy}{dx} = \frac{B - \sqrt{B^2 - 4AC}}{2A} = -\frac{x}{y}.
\]

\begin{verbatim}
обеспечивает  \overline{A} = \overline{C} = 0  ),
\end{verbatim}

\pandocbounded{\includegraphics[keepaspectratio]{images/_page_172_Figure_2.jpeg}}

Фіг. 23.3. Нові координати карактора.

Інтегруючи ці звичайні диференціальні рівняння методом розділення
змінних, ми отримуємо два неявних співвідношення (за бажанням можна
виразити y через x)

\[
y^2 - x^2 = \text{const},
y^2 + x^2 = \text{const}.
\]

Следовательно, новые координаты следует вводить по формулам

\[
\xi = y^2 - x^2
\]

,\\
\(\eta = y^2 + x^2\) .

Эти новые координаты изображены на рис. 23.3. Теперь следует получить
новое уравнение. Вычислим все коэффициенты

\(\widetilde{A} = 0\) (так и должно быть; мы специально выбирали
\(\xi \xi\) и \(\eta\) так, чтобы этот коэффициент обращался в нуль).

\[
\vec{B} = 2A\xi_x\eta_x + B(\xi_x\eta_y + \xi_y\eta_x) + 2C\xi_y\eta_y = -16x^2y^2
\]

,

\(\overline{C} = 0\) (по той же причине, что и \(\overline{A} = 0\) ),

\[
\bar{D} = A\xi_{xx} + B\xi_{xy} + C\xi_{yy} + D\xi_x + E\xi_y = -2(x^2 + y^2)
\widetilde{E} = A\eta_{xx} + B\eta_{xy} + C\eta_{yy} + D\eta_x + E\eta_y = 2(y^2 - x^2),
\tilde{F} = F = 0
\]

,

\[
\overline{G} = G = 0
\]

,

подставим их в уравнение

\[
\overline{A}u_{\xi\xi} + \overline{B}u_{\xi\eta} + \overline{C}u_{\eta\eta} + \overline{D}u + \overline{E}u_{\eta} + \overline{F}u = \overline{G}
\]

и получим

\[
u_{\xi\eta} = \frac{-(x^2 + y^2) u_{\xi} + (y^2 - x^2) u_{\eta}}{8x^2y^3}.
\]

ШАГ 2. Выражая, наконец, x и y через \(\xi\) и \(\eta\) , окончательно
получаем

(23.8)

\[
u_{\xi\eta} = \frac{\eta u_{\xi} - \xi u_{\eta}}{2(\xi^2 - \eta^2)}.
\]

Задача решена. Теперь мы знаем, что нужно сделать, чтобы найти новые
координаты и привести уравнение к канонической форме. \# ЗАУВАЖЕННЯ 1.
На самом деле для гиперболических уравнений существует две канонические
формы. Вторую можно получить из первов заменой переменных вида

\[
\alpha = \alpha (\xi, \eta) = \xi + \eta,
\]

\(\beta = \beta (\xi, \eta) = \xi - \eta.\) Воспользуемся этой заменой
для преобразования уравнения (23.8). Так как

\[
\begin{array}{l} u_{\xi} = u_{\alpha}\alpha_{\xi} + u_{\beta}\beta_{\xi} = u_{\alpha} + u_{\beta}. \\ u_{\eta} = u_{\alpha}\alpha_{\eta} + u_{\beta}\beta_{\eta} = u_{\alpha} - u_{\beta}, \\ u_{\xi\eta} = u_{\alpha\alpha}\alpha_{\eta} + u_{\alpha\beta}\beta_{\eta} + u_{\beta\alpha}\alpha_{\eta} + u_{\beta\beta}\beta_{\eta} = u_{\alpha\alpha} - u_{\beta\beta}, \end{array}
\]

TO

(23 9)

\[
u_{\alpha\alpha} - u_{\beta\beta} = \frac{-\beta u_{\alpha} - \alpha u_{\beta}}{2\alpha\beta}.
\]

При желании можно выразить \(\alpha\) и \(\beta\) через исходные
координаты x и y. В результате получаем

\[
\alpha = \xi + \eta = (y^2 - x^2) + (y^2 + x^2) = 2y^3,
\beta = \xi - \eta = (y^2 - x^2) - (y^2 + x^2) = -2x^2.
\]

\begin{itemize}
\item
  Читач може запитати, чому необхідно класифікувати та зводити їх до
  канонічної форми диференціальних рівнянь у похідних похідних.
\item
  \begin{enumerate}
  \def\labelenumi{\alph{enumi})}
  \tightlist
  \item
    Поділ рівнянь на класи гіперболічних, параболічних та еліптичних
    відповідає поділу фізичних процесів на три основні класи: хвильовий,
    дифузійний і стаціонарний. Математичні особливості розв'язків цих
    трьох типів рівнянь є абсолютно різними.
  \end{enumerate}
\item
  \begin{enumerate}
  \def\labelenumi{\alph{enumi})}
  \setcounter{enumi}{1}
  \tightlist
  \item
    У більшості праць, присвячених розв'язанню гіперболічних задач,
    припускається, що початкове рівняння записується у канонічній формі,
    тобто у вигляді
  \end{enumerate}
\end{itemize}

\[
u_{\xi\xi} - u_{\eta\eta} = \Psi(\xi, \, \eta, \, u_{\xi}, \, u, u_{\eta}).
\]

Якщо у нас є рівняння і ми хочемо отримати його розв'язок, ми повинні
перетворити його у канонічну форму і використати відомі результати.

\begin{enumerate}
\def\labelenumi{\alph{enumi})}
\setcounter{enumi}{2}
\tightlist
\item
  Для чисельного розв'язання канонічного гіперболічного рівняння було
  створено багато різних комп'ютерних програм.
\end{enumerate}

\begin{verbatim}
Функція  \Phi  (  \xi  ,  \xi \eta  , u,  u_{\xi}  ,  u_{\eta}  ) вводиться у VM  \Theta  як підпрограма, тому рівняння спочатку має бути канонічним. Як тільки розв'язок знайдено в нових координатах, його завжди можна перетворити назад на початкові координати.
\end{verbatim}

\section{ЗАВДАННЯ}\label{ux437ux430ux432ux434ux430ux43dux43dux44f-1}

\begin{enumerate}
\def\labelenumi{\arabic{enumi}.}
\tightlist
\item
  Визначте, чи є наступні рівняння гіперболічними, параболічними чи
  еліптичними:
\end{enumerate}

\begin{verbatim}
- a) u_{xx}u_{xy}=0,
   - б) u_{tt}=u_{xx}+u_x+hu,
   - в) u_{xx}+3u_{yy}=\sin x,
   - г) u_{xx}+u_{yy}=f(x,y),
   - д) u_{rr}+\frac{1}{r}u_r+\frac{1}{r^2}u_{\theta\theta}=f(r,\theta).
\end{verbatim}

\begin{enumerate}
\def\labelenumi{\arabic{enumi}.}
\setcounter{enumi}{1}
\item
  Отримайте співвідношення (23.4), (23.5) та (23.6).
\item
  Переконайтеся, що рівняння
\end{enumerate}

\[
3u_{xx} + 7u_{xy} + 2u_{yy} = 0
\]

є гіперболічним при всіх x та y, і знайдіть його характеристики.

\begin{enumerate}
\def\labelenumi{\arabic{enumi}.}
\setcounter{enumi}{3}
\tightlist
\item
  Зведіть рівняння з попередньої задачі до канонічної форми
\end{enumerate}

\[
u_{\xi\eta} = \Phi(\xi, \eta, u, u_{\xi}, u_{\eta}).
\]

\begin{enumerate}
\def\labelenumi{\arabic{enumi}.}
\setcounter{enumi}{4}
\tightlist
\item
  Продовжте задачу 4 і приведіть це рівняння до другої канонічної форми
\end{enumerate}

\[
u_{\alpha\alpha} - u_{\beta\beta} = \Psi(\alpha, \beta, u, u_{\alpha}, u_{\beta}).
\]

\begin{enumerate}
\def\labelenumi{\arabic{enumi}.}
\setcounter{enumi}{5}
\tightlist
\item
  Найдите характеристики уравнения
\end{enumerate}

\[
u_{xx} + 4u_{xy} = 0.
\]

Перетворіть рівняння на нові координати, розв'яжіть його і поверніться
до координат x і y, щоб отримати розв'язок початкової здогадки.

\section{ХВИЛЬОВЕ РІВНЯННЯ У ВІЛЬНОМУ ПРОСТОРІ (ДВОВИМІРНІ ТА ТРИВИМІРНІ
ЗАДАЧІ)}\label{ux445ux432ux438ux43bux44cux43eux432ux435-ux440ux456ux432ux43dux44fux43dux43dux44f-ux443-ux432ux456ux43bux44cux43dux43eux43cux443-ux43fux440ux43eux441ux442ux43eux440ux456-ux434ux432ux43eux432ux438ux43cux456ux440ux43dux456-ux442ux430-ux442ux440ux438ux432ux438ux43cux456ux440ux43dux456-ux437ux430ux434ux430ux447ux456}

МЕТАЛЕКЦІЇ: Розв'яжіть задачу з початковими умовами

(УЧП)

\[
u_{tt} = e^{2} (u_{xx} + u_{yy} + u_{zz}),
\begin{cases} -\infty < x < \infty, \\ -\infty < y < \infty, \\ -\infty < z < \infty, \end{cases}
\]

(НУ)

\[
\begin{cases} u(x, y, z, 0) = \varphi(x, y, z), \\ u_{t}(x, y, z, 0) = \psi(x, y, z) \end{cases}
\]

в трехмерном пространстве и показать, что полученное решение
удовлетворяет \emph{принципу Гюйгенса}. Для решения соответствующей
двумерной задачи

(УЧП)

\[
u_{tt} = c^2 (u_{xx} + u_{yy}), \quad
\begin{cases} -\infty < x < \infty, \\ -\infty < y < \infty, \end{cases}
\]

(НУ)

MATHBLK0090X

використовується метод спуску, і доведено, що двовимірні розв'язки не
задовольняють принцип Гюйгенса. І нарешті, використовуючи метод спуску,
показано, що розв'язок одномірного рівняння виражається формулою
д'Аламбера, з якою ми познайомилися раніше.

У попередніх лекціях ми розглядали задачу з початковими умовами для
нескінченної вібруючої струни і показали, що її розв'язок задається
формулою д'Аламбера. Читачеві має бути зрозуміло, що в тривимірному
просторі одномірне хвильове рівняння описує плоскі хвилі. Наприклад,
звукові або електромагнітні хвилі на достатньо великих відстанях від
джерел можна вважати плоскими хвилями і, отже, описуватися одномірним
рівнянням. Існує така термінологія:

\begin{itemize}
\tightlist
\item
  \begin{enumerate}
  \def\labelenumi{\arabic{enumi})}
  \tightlist
  \item
    одномірні хвилі називаються плоскими хвилями,
  \end{enumerate}
\item
  \begin{enumerate}
  \def\labelenumi{\arabic{enumi})}
  \setcounter{enumi}{1}
  \tightlist
  \item
    Двовимірні хвилі називають циліндричними хвилями.
  \end{enumerate}
\end{itemize}

\begin{enumerate}
\def\labelenumi{\arabic{enumi})}
\setcounter{enumi}{2}
\tightlist
\item
  тривимірні хвилі називаються сферичними хвилями.
\end{enumerate}

Іншими словами, одномірне хвильове рівняння може описувати або плоскі
хвилі в просторах більшої кількості розмірів, або коливання струн.
Завдання лекції --- узагальнити формулу д'Аламбера для випадку двох і
трьох вимірів.

\section{Хвилі у тривимірному
просторі}\label{ux445ux432ux438ux43bux456-ux443-ux442ux440ux438ux432ux438ux43cux456ux440ux43dux43eux43cux443-ux43fux440ux43eux441ux442ux43eux440ux456}

Розглянемо сферичні хвилі у тривимірному просторі за заданих початкових
умов, тобто розв'яжемо задачу з початковими умовами

(24.1)

\[
\begin{cases} -\infty < x < \infty, \\ -\infty < y < \infty, \\ -\infty < z < \infty, \end{cases}
\]

(24.1)

\[
\begin{cases} u(x, y, z, 0) = \varphi(x, y, z), \\ u_{t}(x, y, z, 0) = \psi(x, y, z). \end{cases}
P_{2}
\begin{aligned} u_{tt} &= c^2 \Delta u, \\ (HY) & \begin{cases} u(x, y, z, 0) = 0, \\ u_t(x, y, z, 0) = \psi(x, y, z), \end{cases} \end{aligned}
\]

\begin{verbatim}
где экаком  \Delta  обозначен дифференциальный оператор
\end{verbatim}

Цю задачу можна розв'язати методом \(\Phi\) Юр'єра \(^{1}\) ), і
виявляється, що розв'язок можна представити у вигляді

\[
P_{1}
\]

где \(\overline{\psi}\) --- \emph{среднее} значение начального
распределения \(\psi\) \emph{по сфере} радиуса ct с центром в точке (x,
y, z), т. е.

\[
+ ct \cos \varphi  )  (ct)^2 \sin \varphi d\varphi d\theta  .
\]

\begin{verbatim}
Аргументи функції  \psi  проходити по поверхні сфери, якщо значення  \theta  і  \phi  змінюються відповідно в межах  [0, 2\pi]  і  [0, \pi]  (рисунок 24.1). Цю формулу можна інтерпретувати так:
\end{verbatim}

См Владимиров В. С. Уравнения математической физики. --- М.: Наука, 1971
и др. издания. --- Прим ред.

дующим образом: каждая точка пространства излучает расходящуюся (со
скоростью c) сферическую волну, через t секунд точка с координатами (x,
y, z) окажется под воздействием начального Р И С

\pandocbounded{\includegraphics[keepaspectratio]{images/_page_177_Picture_3.jpeg}}

. 24 1. Решение является средним значением начатьного распределения на
сфере.

возмущения, сосредоточенного на сфере раднуса ct с центром в заданной
точке (см. рис. 24.2).

Формула (24.3) дозволяє обчислити розв'язок задачі для більшості
початкових умов на комп'ютері. Можливо, читачу буде цікаво знайти
розв'язок для деяких простих функцій f.

\pandocbounded{\includegraphics[keepaspectratio]{images/_page_177_Picture_7.jpeg}}

\begin{verbatim}
Фіг. 24.2. Початкове збурення  \psi  поширюється з кожної точки у всіх напрямках.
\end{verbatim}

Щоб отримати повне рішення, розглянемо другу половину проблеми

(24.4)

\[
\begin{aligned} (\text{YAII}) & u_{tt} = c^2 \Delta u, \\ (\text{HY}) & \begin{cases} u = 0, \\ u_t = \varphi, \end{cases} \end{aligned}
u_h(r, \theta) = \sum_{n=0}^{\infty} r^n [a_n \cos(n\theta) + b_n \sin(n\theta)].
\Delta u = -\frac{r^2}{2} - \frac{r^3}{2} \cos{(2\theta)}
u = \frac{\partial}{\partial t} [t\overline{\varphi}].
u_p(r, \theta) = Ar^4 + Br^4 \cos(2\theta).
u(x, t) = \frac{1}{2c} \int_{x-ct}^{x+ct} \varphi(s) ds.
u_p(r, \theta) = -\frac{r^4}{32} - \frac{r^4}{24} \cos(2\theta).
u_t(x, t) = \frac{1}{2} [\varphi(x + ct) + \varphi(x - ct)].
\sum_{n=0}^{\infty} [a_n \cos(n\theta) + b_n \sin(n\theta)] - \frac{1}{32} - \frac{1}{24} \cos(2\theta) = 0.
(\text{УЧП}) \quad u_{tt} = c^2 \Delta u, \quad (x, y, z) \in \mathbb{R}^3,
u_1(r, \theta) = \frac{1}{32} + \frac{1}{24}r^2\cos(2\theta) - \frac{r^4}{24}\cos(2\theta) - \frac{r^4}{32} = \frac{(r^4 - 1)}{32} - \frac{(r^4 - r^2)}{24}\cos(2\theta).
(\text{HV}) \quad
u = u_0 + u_1 = r \cos \theta - \frac{(r^4 - 1)}{32} - \frac{(r^4 - r^2)}{24} \cos (2\theta).
\begin{cases} u = \varphi, \\ u_t = \psi \end{cases}
P_{2}
\Delta u = 0, \quad 0 < r < 1 + \frac{1}{4} \sin \theta,
\]

записывается в виде

\[
u \left( 1 + \frac{1}{4} \sin \theta, \ \theta \right) = \cos \theta, \quad 0 \le \theta \le 2\pi.
\]

где \(\overline{\phi}\) и \(\overline{\psi}\) --- средние значения
функций \(\phi\) и \(\psi\) на сфере раднусом ct с центром в точке (x,
y, z).

Цей розв'язок зазвичай називають формулою Пуассона для хвильового
рівняння у тривимірному просторі. Це природне узагальнення формули
д'Аламбера на тривимірний випадок. Найважливішим у формулі Пуассона є
те, що обидва інтеграли в \(\bar{\phi}\) і \(\bar{\psi}\) взяті з
поверхні сфери. Цей факт дозволяє нам надати таку важливу інтерпретацію
рішення.

У момент \(t=t_1\) розв'язок u у (x,y,z) залежить лише від початкових
розподілів сфери з радіусом \(ct_1\) центрованим на (x,y,z) (рисунок
24.3). Тепер припустимо, що початкові розподіли \(\phi\) і \(\psi\)
дорівнюють нулю всюди, крім малої сфери (див. рисунок 24.3). З часом
радіус P И С

\pandocbounded{\includegraphics[keepaspectratio]{images/_page_179_Picture_3.jpeg}}

. 24.3. Схема поясияет, как начальное возмущение влияет на точку (к, ч,
г).

Сфера, що оточує точку (x, y, z), зростає зі швидкістю c, так що через
\(t_{\hat{z}}\) секунди ця сфера починає перетинати область початкового
збурення, і тому розв'язок \(u(x, y, z, t)\) стає ненульовим. У
\(t_{\hat{z}} < t < t_{\hat{a}}\) розв'язок U у (x, y, z) залишається
ненульовим, оскільки сфера продовжує перетинати область початкового
збурення, але на \(t=t_{\hat{a}}\) розв'язок знову стає нульовим. Іншими
словами, розв'язок, що поширюється з початкової області збурення, має
чітко визначену задню кромку. Цей загальний принцип відомий як принцип
Гюйгенса для тривимірного простору. Виявляється, що передній край хвилі
завжди чітко окреслений, але задній край чітко окреслений лише в
просторах розмірності 3, 5, 7\ldots. Ми вже знаємо з формули д'Аламбера,
що початкові умови Р и

\pandocbounded{\includegraphics[keepaspectratio]{images/_page_185_Figure_3.jpeg}}

с. 25.1. Графики функции \(f(x)\) = 1 и ее преобразования.

\section{Свойства
преобразований}\label{ux441ux432ux43eux439ux441ux442ux432ux430-ux43fux440ux435ux43eux431ux440ux430ux437ux43eux432ux430ux43dux438ux439}

Прежде чем приступить к решению задач, мы должны получить некоторые
свойства этих преобразований.

Якщо you(x, t) --- функція \(\partial \theta yx\) змінних, і ми
виконуємо перетворення змінної x, то

(действительная форма отображения w=z\^{}2 ).

Ця форма нам знадобиться пізніше. У цьому прикладі це цікаво, оскільки
показує, як гіперболи в площині z переходять до координатних прямих u =
const і v = const у площині w.

Тепер перейдемо до вивчення одного конкретного типу функцій комплексної
змінної, які виконують конформні відображення.

\section{Определение конформного
отображения}\label{ux43eux43fux440ux435ux434ux435ux43bux435ux43dux438ux435-ux43aux43eux43dux444ux43eux440ux43cux43dux43eux433ux43e-ux43eux442ux43eux431ux440ux430ux436ux435ux43dux438ux44f}

\begin{verbatim}
Відображення w=f(z) комплексної площини z у комплексну площину w називається конформним у точці z_{\bullet} площини z, якщо похідна f'(z_{\bullet}) \neq 0 . Відображення f(z) називається конформним у регіоні D, якщо f'(z) \neq 0 в кожній точці регіону D.
\end{verbatim}

\begin{verbatim}
У застосуванні f(z) = z^2 конформна з усіх речей, окрім точки z = 0, z'(z) = 2z \neq 0 для всіх z \neq 0 . З іншого боку, якість e^z відповідає всім z-областям, f'(z) = e^z \neq 0 є конформною. Яка мета конформних стосунків? Відповід полягає полягає в тому, що розвиток рівня Лапласа \phi_{xx} + \phi_{yy} = 0 у співаючій області області змін x i y, ми можемо розглядати цю область як площу комплексної зміни z. Відожажуючи w = f(z \phi_{xy} + \phi_{uy} = 0 ), рівень Лапласа \varphi_{xx} + \varphi_{yy} = 0 перетворюватися на новому диференціальному рівні частково, як у нових координатах і v \varphi_{nn} + \varphi_{nn} = 0 \phi_{xx} + \phi_{uu} = 0 . Водночас я придумав ідею знайти конформну межу, як перенести область з комплексною межею та область з простою межею (припускають, що область є w=z^* перенести межу першого квадранта площі z на дійсну всієї площі w).
\end{verbatim}

Ось кілька прикладів, які покажуть, як можна розв'язати диференціальні
рівняння у похідних за допомогою конформних відображень.

Уравнение Лапласа в верхней полуплоскости.

Припустимо, що нам потрібно розв'язати наступну задачу Діріхле у верхній
половині видимості (рис. 47.3):

(47.1)

(Зверніть увагу, що перетворення було виконано на змінній x, і отримана
послідовність залежить лише від часу t.)

А как быть с производными? Ниже приводятся несколько полезных формул

\[
\begin{aligned} (Y & \Pi) & \varphi_{xx} + \varphi_{yy} = 0, & -\infty < x < \infty, & 0 < y < \infty, \\ (47.1) & \varphi(x, 0) = \begin{cases} 0, & |x| > 1, \\ 1, & |x| \le 1. \end{cases} \end{aligned}
w = \ln\left\{\frac{z-1}{z+1}\right\}
\begin{split} S\left[u_{t}\right] &= \frac{dS\left[u\right]}{dt}, \quad S\left[u_{tt}\right] = \frac{d^{2}S\left[u\right]}{dt^{2}}, \\ S\left[u_{xx}\right] &= -\left[n\pi/L\right]^{2}S\left[u\right] + \frac{2n\pi}{L^{2}}\left[u\left(0,\ t\right) + (-1)^{n+2}u\left(L,\ t\right)\right], \\ C\left[u_{xx}\right] &= -\left[n\pi/L\right]^{2}C\left[u\right] - \frac{2}{L}\left[u_{x}\left(0,\ t\right) + (-1)^{n+1}u_{x}\left(L,\ t\right)\right]. \end{split}
\]

\pandocbounded{\includegraphics[keepaspectratio]{images/_page_189_Picture_10.jpeg}}

Фіг. 26.1. Основна ідея принципу суперпоемації; L --- диференціальний
оператор; f --- вхід; У---виходь.

Ця базова ідея може використовуватися для розв'язання складних змішаних
задач, розбиваючи їх на прості підзавдання, які розв'язуються за
допомогою

кожне з них, а потім підсумовування отриманих результатів (звісно,
диференціальне рівняння і граничні умови мають бути лінійними).

\section{Розкладання змішаної задачі на дві
простіші}\label{ux440ux43eux437ux43aux43bux430ux434ux430ux43dux43dux44f-ux437ux43cux456ux448ux430ux43dux43eux457-ux437ux430ux434ux430ux447ux456-ux43dux430-ux434ux432ux456-ux43fux440ux43eux441ux442ux456ux448ux456}

Предположим, нам необходимо решить линейную задачу (обозначим ее P)

(УЧП)

\begin{enumerate}
\def\labelenumi{\arabic{enumi}.}
\setcounter{enumi}{17}
\tightlist
\item
\end{enumerate}

(Р) (ГУ)

\[
\begin{cases} 1-|x|, |x| < 1 \\ 0, |x| > 1 \end{cases}
\begin{cases} u(0, t) = 0, \\ u(1, t) = 0, \\ u(x, 0) = \sin(2\pi x), \quad 0 \le x \le 1. \end{cases}
\cos(ax)
\frac{2}{\pi} \frac{\sin a\omega}{\omega}
\]

Перед нами неоднородное уравнение теплопроводности, так что метод
разделения переменных неприменим. Можпо, конечно, воспользоваться
конечным синус-преобразованием по переменной x или преобразованием
Лапласа по переменной t, но вместо этого мы рассмотрим две подзадачи

\[
\frac{2}{\pi} \frac{\sin a\omega}{(\omega^2+a^2)^2}
\frac{\pi}{2} e^{-a+\omega}
\begin{cases} u(0, t) = 0, \\ u(1, t) = 0, \\ u(1, t) = 0, \end{cases}
\frac{\pi}{2} e^{-a+\omega}
\frac{\pi}{2} e^{-a+\omega}
\frac{\pi}{2} e^{-a+\omega}
\frac{\pi}{2} e^{-a+\omega}
\]

Каждая из этих задач легко решается и в то же время ясно, что
\(\mathit{suma}\) решений задач \(P_1\) и \(P_2\) является решением
исходной вадачи, т. е.

\[
\frac{2ax}{(x^2+a^2)^2}
\]

Решение задачи \(P_{i}\) Розв'язання задачі \(P_i\) Розв'язання задачі
\(P_i\) Можна показати, що в загальному випадку розв'язок задачі

\[
\frac{\cos(ax), |x| > \pi/2a}{0, |x| > \pi/2a}
\frac{1-|x|, |x| < 1}{0, |x| > 1}
\begin{cases} u_t = u_{xx} + f(x, t), \\ u(0, t) = 0, \\ u(1, t) = 0, \\ u(x, 0) = \varphi(x) \end{cases}
\cos(ax)
\cos(ax)
\frac{\pi}{2} [\delta(\omega+a) + \delta(\omega-a)]
\begin{cases} u_t = u_{xx} + f(x, t), \\ u(0, t) = 0, \\ u(1, t) = 0, \\ u(x, 0) = 0 \end{cases}
f(x) = \int_{0}^{\infty} F(\omega) \sin(\omega x) d\omega
F(\omega) = \frac{?}{\pi} \int_{0}^{\infty} f(x) \sin(\omega x) dx
0 < x < \infty
\begin{cases} u_t = u_{xx}, \\ u(0, t) = 0, \\ u(1, t) = 0, \\ u(x, 0) = \varphi(x). \end{cases}
0 < \omega < x
\]

f'\,'(0)

\[
-\omega^2 F(\omega) + \frac{2}{\pi} \omega f(0)
\frac{1}{\alpha} F\left(\frac{\omega}{\alpha}\right)
\]

проследим, как оно решается методом конечного синус-преобразования. Вот
что мы должны сделать: разложить еходное воздействие \(f(x, t)\) на
компоненты, найти отклик \(U_n\) на каждую компоненту, а затем сложить
эти отклики. Математически это может быть и не совсем очевидно, но
давайте внимательно проследим за ходом решения. Разложим входящие в
уравнение с частными производными члены в ряды Фурье по синусам:

\[
\frac{2\omega}{\pi (\alpha^2 + \omega^2)}
x^{-1/2}
[2/\pi\omega]^{1/2}
\]

where

\[
\frac{2}{\pi\omega} [1-\cos(\omega a)]
\]

1

\[
e^{-1}a + \omega
\]

Зверніть увагу, по-перше, що коефіцієнти \(A_n\) , \(B_n\) \(F_n\)
дійсно є функціями часу t, оскільки ми застосували синусоїдальне
перетворення до функції двох змінних x, t, u, а по-друге, ми розклали
вхідний ефект \(f(x,t)\) на прості компоненти \(F_n(t)\) . Тепер
знайдіть відповіді \(U_n(t)\) для кожного компонента \(F_n(t)\) , потім
додайте всі відповіді \(U_n(t)\) і отримайте розв'язок \(u(x,t)\).

Щоб знайти відповіді \(U_n(t)\) , необхідно робити невеликі розрахунки з
коефіцієнтами \(A_n(t)\) і \(B_n(t)\) , щоб у вирази інтегрування була
включена лише функція u, а не її похідні \(u_i\) і \(u_{xx}\) .
Інтегруючи на частини, ми отримуємо

\[
\frac{1}{2}e^{-\alpha}\sin\omega
\frac{1-e^{-a\omega}}{\omega}
\begin{split} A_n(t) &= 2\int_0^1 u_t \sin(n\pi x) \, dx = \frac{d}{dt} \bigg[ 2\int_0^1 u(x, t) \sin(n\pi x) \, dx \bigg] = \frac{dU_n(t)}{dt} \,, \\ B_n(t) &= 2\int_0^1 u_{xx} \sin(n\pi x) \, dx = -(n\pi)^2 \, U_n(t) \,+ \\ &\quad + 2n\pi \, \big[ u(0, t) + (-1)^{n+1} \, u(1, t) \big], \end{split}
\frac{2}{\pi}F''(\omega)
f(x) = \int_{0}^{\infty} F(\omega) \cos(\omega x) d\omega \qquad \qquad F(x) = \frac{2}{\pi} \int_{0}^{\infty} f(x) \cos(\omega x) dx
0 < x < \infty \qquad \qquad 0 < \omega < \infty
\begin{cases} u(0, t) = 0, \\ u(1, t) = 0 \end{cases}
\]

j'\,'(x)

\[
S_n = \frac{2}{\pi} \int_0^{\pi} f(x) \sin(nx) dx.
\]

в последнее соотношение, получаем

\[
y = \frac{\pi}{b - a} (x - a).
\]

и, следовательно, (26.1) принимает вид

\[
f(x) = \sum_{n=1}^{\infty} S_n \sin(nx).
\]

Поскольку это разложение является тождеством по x, его коэффициенты
должны быть равны нулю, \(\tau\) . е.

\[
f(x) = \sum_{n=1}^{\infty} S_n \sin(nx)
\]

.

Таким образом, мы получили связь между входом и выходом системы, между
коэффициентами \(F_n\) и \(U_n\) . Прежде чем решать уравнения для
\(U_n(t)\) , обратимся к начальному условию

\[
S_n = \frac{2}{\pi} \int_{0}^{\pi} f(x) \sin(nx) dx
\]

Если разложить \(u(x, 0)\) в ряд Фурье по синусам и приравнять его нулю:

\[
0 \le x \le \pi
\]

то определяются начальные условия для функции \(U_n\left(t\right)\) , т.
е.

\[
n = 1, 2, ...
\]

, \(n = 1, 2, \dots\) Мы получили разложение нашей смешанной задачи на
ряд простых задач об отклике системы на элементарные входные
воздействия:

(OДУ)

f'\,'(x)

(HУ) \(U_n(0) = 0, n = 1, 2, ...\) Каждое из этих уравнений легко
решается методом интегрирующего множителя (или с помощью преобразования
Лапласа). В любом случае получаем

\[
\frac{\pi}{2} e^{ax}
\]

Таким образом, мы получили отклики \(U_n(t)\) на простейшие входные
воздействия \(F_n(t)\) . Для получения решения исходной задачи осталось
совершить последний шаг---сложить все отклики

\[
C_n = \frac{2}{\pi} \int_0^{\pi} f(x) \cos(nx) dx.
\]

Читатель, должно быть, заметил, что каждый отклик \(U_n(t)\) умножается
на весовой множитель \(\sin(n\pi x)\) . Конечно, в разложении функции
\(f(x, t)\) на компоненты \(F_n(t)\) присутствуют те же весовые
множители.

\section{ЗАМЕЧАНИЕ}\label{ux437ux430ux43cux435ux447ux430ux43dux438ux435}

Використання скінченного синусоїдального перетворення дозволяє записати
розв'язок у вигляді ряду, але для більшості інших інтегральних
перетворень розв'язки представлені як інтеграли (неперервний розклад).
Читача можна попросити дати інтерпретацію проблеми

\[
y = \frac{\pi}{b-a}(x-a).
f(x) = \frac{C_0}{2} + \sum_{n=1}^{\infty} C_n \cos(nx).
f(x) = \frac{C_0}{2} + \sum_{n=1}^{\infty} C_n \cos(nx) \qquad C_n = \frac{2}{\pi} \int_{0}^{\pi} f(x) \cos nx \, dx
0 \le x \le \pi \qquad n = 0, 1, 2, \dots
\]

\begin{verbatim}
\begin{array}{llll} (\mathrm{Y} + \Pi) & u_{tt} = u_{xx} + \sin{(3\pi x)}, & 0 < x < 1, & 0 < t < \infty, \\ (\mathrm{\Gamma} \mathrm{Y}) & \left\{ \begin{array}{ll} u\left(0, \ t\right) = 0, \\ u\left(1, \ t\right) = 0, \\ \end{array} \end{array}
\end{verbatim}

f'\,'(x)

\[
-n^{2}C_{n} - \frac{2}{\pi} \{f'(0) - (-1)^{n} f'(n)\}
\begin{array}{ll} u\left(x, \ 0\right) = \sin{(\pi x)}, \\ u_{t}\left(x, \ 0\right) = 0, \end{array}
\]

Для каждой из подзадач выберите метод решения по своему BKYCY,

\begin{enumerate}
\def\labelenumi{\arabic{enumi}.}
\setcounter{enumi}{2}
\tightlist
\item
  Нехай \(u_1\) і \(u_2\) є розв'язками наступних рівнянь:
\end{enumerate}

\begin{enumerate}
\def\labelenumi{\alph{enumi})}
\tightlist
\item
\end{enumerate}

;

\[
b)  u_t = u_{xx} + e^t  ;
\]

\begin{enumerate}
\def\labelenumi{\Alph{enumi})}
\setcounter{enumi}{1}
\tightlist
\item
\end{enumerate}

\begin{enumerate}
\def\labelenumi{\arabic{enumi}.}
\setcounter{enumi}{4}
\tightlist
\item
  \(\cos(mx), m=1, 2, \dots\)
\end{enumerate}

;

\[
F)  u_t = u_{xx} + u^2  .
\]

\(\mathbf{\Pi}\) для яких рівнянь функція \(u_1 + u_2\) також буде
розв'язком? Які висновки ви можете зробити зі своїх відповідей?

\begin{enumerate}
\def\labelenumi{\arabic{enumi}.}
\setcounter{enumi}{3}
\tightlist
\item
  Найдите четыре смещанные задачи, сумма решений которых дает решение
  следующей задачи:
\end{enumerate}

\[
\begin{aligned} u_t &= u_{xx} + f(x, t), & 0 < x < 1, & 0 < t < \infty, \\ & \begin{cases} u(0, t) = g_1(t), \\ u(1, t) = g_2(t), \end{cases} & 0 < t < \infty, \\ & u(x, 0) = \varphi(x), & 0 \le x \le 1. \end{aligned}
\]

\begin{enumerate}
\def\labelenumi{\arabic{enumi}.}
\setcounter{enumi}{4}
\tightlist
\item
  Решите задачу Коши
\end{enumerate}

(OAY)

x

(HY) \(U_n(0) = 0.\) Можете ли вы проверить полученное решение?

\begin{enumerate}
\def\labelenumi{\arabic{enumi}.}
\setcounter{enumi}{5}
\tightlist
\item
  Предположим, что функции \(u_1\) и \(u_2\) удовлетворяют линейным
  однородным граничным условиям
\end{enumerate}

\begin{verbatim}
\begin{cases} u_x(1, t) + h_2 u(1, t) = 0.  Чи відповідає функція u_i + u_i цим умовам? , .
\end{verbatim}

Лекция 27

\section{РІВНЯННЯ ПЕРШОГО ПОРЯДКУ (ХАРАКТЕРИСТИЧНИЙ
МЕТОД)}\label{ux440ux456ux432ux43dux44fux43dux43dux44f-ux43fux435ux440ux448ux43eux433ux43e-ux43fux43eux440ux44fux434ux43aux443-ux445ux430ux440ux430ux43aux442ux435ux440ux438ux441ux442ux438ux447ux43dux438ux439-ux43cux435ux442ux43eux434}

МЕТАЛЕКЦІЇ: Введіть поняття рівняння з похідними в похідних першого
порядку (поки що ми розглядали рівняння другого порядку) і ознайомтеся з
важливим методом розв'язання задач із початковими умовами --- методом
характеристик. Проблема, яку ми зараз розглянемо, виглядає так:

(УЧП)

\begin{verbatim}
| _                                |,   -\infty < x < \infty  ,   0 < t < \infty  ,
\end{verbatim}

\[
(HV)   u(x, 0) = \varphi(x), \quad -\infty < x < \infty.
\]

\begin{verbatim}
Слід одразу зазначити, що вперше ми розв'яжемо задачу з змінними коефіцієнтами. Виявляється, що якщо ми рухаємося з координат (x, t) на нові, характеристичні координати (s, \tau), то наша UFP перетвориться на звичайне диференціальне рівняння. Тоді можна розв'язати звичайне диференціальне рівняння, тобто знайти u(s, \tau) , і на останньому кроці виразити s і t за x і t і отримати u(x, t).
\end{verbatim}

Напомним читателю, что, когда мы решали уравнение дифф

\[
\Delta u = u_{xx} + u_{yy},   -\infty < x < \infty  ,   0 < t < \infty
\]

\begin{verbatim}
Константа \alpha^2 виконувала роль коефіцієнта дифузії, а v — швидкість середовища. Очевидно, якщо \alpha=0 (тобто дифузії немає), то розв'язок буде зсуватися як ціле число вздовж осі x зі швидкістю v (залишається лише конвекція). Іншими словами, якщо початкова умова — u(x,0)=\varphi(x) , то відповідне розв'язання рівняння
\end{verbatim}

\[
\Delta u = u_{rr} + \frac{1}{r}u_r + \frac{1}{r^2}u_{\theta\theta}
\]

\begin{verbatim}
будет иметь вид   u(x, t) = \varphi(x-vt)  .
\end{verbatim}

Це дає нам підстави вважати, що розв'язок рівняння першого порядку

\[
\Delta u = u_{xx} + u_{yy} + u_{zz}
\]

переносится вдоль потока со скоростью

\[
\Delta u = u_{rr} + \frac{1}{r}u_r + \frac{1}{r^2}u_{\theta\theta} + u_{zz}.
\]

Звісно, якщо a і b --- константи, то розв'язок --- це хвиля, що
рухається з постійною швидкістю. Якщо \(a(x, t)\) та \(b(x, t)\)

\pandocbounded{\includegraphics[keepaspectratio]{images/_page_196_Figure_10.jpeg}}

\begin{verbatim}
Фіг. 27.1. Початкове значення в точці x впливає на розв'язок лише вздовж характеристики у площині змінних (x,t). Параметр s змінюється вздовж характеристик від нуля до нескінченності, починаючи з початкової лінії умови;   \tau — сталий за кожною характеристикою;   \tau змінюється вздовж початкової лінії стану.
\end{verbatim}

залежать від x і t, то швидкість поширення змінюється і вздовж потоку, і
в часі, тому початкова крива може істотно деформуватися. Тут справді
виникає пряма аналогія з конвекцією.

Розв'язок рівняння першого порядку спирається на той фізичний факт, що
початкове збурення в точці x переноситься у площину (x,t) уздовж прямої,
яку називають характеристикою (див. рис. 27.1).

Таким чином, хвильова задача істотно відрізняється від рівняння
теплопровідності \(u_t = u_{xx}\),для якого початкове значення в точці x
впливає на розв'язок у всіх точках простору і в будь-який момент часу.

\pandocbounded{\includegraphics[keepaspectratio]{images/_page_197_Figure_5.jpeg}}

Фіг. 27.2. Початкове збурення \(u(x,0)\) у точці породжує дві хвилі.
Збурення струни, згідно з хвильовим рівнянням, поширюється вздовж двох
характеристик.

Початкове збурення струни в точці x впливає на розв'язок уздовж двох
характеристичних ліній у площині (x,t), що відповідають двом хвилям, які
рухаються в протилежних напрямках (див. рис. 27.2).

\begin{verbatim}
Природно ввести дві нові координати s і \tau (замість x та t), щоб
\end{verbatim}

\begin{itemize}
\tightlist
\item
  s змінювалася вздовж характеристичних кривих,
\end{itemize}

\begin{verbatim}
- \tau параметризувала початкову криву, бажано вздовж лінії t=0.
\end{verbatim}

Спочатку розглянемо нову координату s. Якщо s вибрано так, щоб
задовольнити наведену умову, то рівняння

\[
\Delta u = u_{rr} + \frac{2}{r}u_r + \frac{1}{r^2}u_{\theta\theta} + \frac{\cot\theta}{r^2}u_0 + \frac{1}{r^2\sin\theta}u_{\phi\phi}
\]

превратится в обыкновенное дифференциальное уравнение К о н е О д н

6 Використовуючи формули, наведені в тексті, розв'яжіть систему
алгебраїчних рівнянь (39.5).

Лекция 40

\section{ПОРІВНЯННЯ АНАЛІТИЧНИХ РІШЕНЬ ІЗ ЧИСЕЛЬНИМИ МЕТАЛЕКЦІЯМИ:
Порівняння переваг і недоліків аналітичних і чисельних методів
розв'язання диференціальних рівнянь у часткових похідних. Важлива
математична задача ідентифікації фізичних величин (параметрична
ідентифікація) наведена. У зв'язку з цим завданням розглядається один
важливий приклад з
біології.}\label{ux43fux43eux440ux456ux432ux43dux44fux43dux43dux44f-ux430ux43dux430ux43bux456ux442ux438ux447ux43dux438ux445-ux440ux456ux448ux435ux43dux44c-ux456ux437-ux447ux438ux441ux435ux43bux44cux43dux438ux43cux438-ux43cux435ux442ux430ux43bux435ux43aux446ux456ux44fux43cux438-ux43fux43eux440ux456ux432ux43dux44fux43dux43dux44f-ux43fux435ux440ux435ux432ux430ux433-ux456-ux43dux435ux434ux43eux43bux456ux43aux456ux432-ux430ux43dux430ux43bux456ux442ux438ux447ux43dux438ux445-ux456-ux447ux438ux441ux435ux43bux44cux43dux438ux445-ux43cux435ux442ux43eux434ux456ux432-ux440ux43eux437ux432ux44fux437ux430ux43dux43dux44f-ux434ux438ux444ux435ux440ux435ux43dux446ux456ux430ux43bux44cux43dux438ux445-ux440ux456ux432ux43dux44fux43dux44c-ux443-ux447ux430ux441ux442ux43aux43eux432ux438ux445-ux43fux43eux445ux456ux434ux43dux438ux445.-ux432ux430ux436ux43bux438ux432ux430-ux43cux430ux442ux435ux43cux430ux442ux438ux447ux43dux430-ux437ux430ux434ux430ux447ux430-ux456ux434ux435ux43dux442ux438ux444ux456ux43aux430ux446ux456ux457-ux444ux456ux437ux438ux447ux43dux438ux445-ux432ux435ux43bux438ux447ux438ux43d-ux43fux430ux440ux430ux43cux435ux442ux440ux438ux447ux43dux430-ux456ux434ux435ux43dux442ux438ux444ux456ux43aux430ux446ux456ux44f-ux43dux430ux432ux435ux434ux435ux43dux430.-ux443-ux437ux432ux44fux437ux43aux443-ux437-ux446ux438ux43c-ux437ux430ux432ux434ux430ux43dux43dux44fux43c-ux440ux43eux437ux433ux43bux44fux434ux430ux454ux442ux44cux441ux44f-ux43eux434ux438ux43d-ux432ux430ux436ux43bux438ux432ux438ux439-ux43fux440ux438ux43aux43bux430ux434-ux437-ux431ux456ux43eux43bux43eux433ux456ux457.}

Можливо, настав час обговорити переваги та недоліки аналітичних і
чисельних розв'язків диференціальних рівнянь у похідних похідних. Спершу
давайте з'ясуємо, що ми маємо на увазі під цими двома типами рішень.

\section{Аналитические
решения}\label{ux430ux43dux430ux43bux438ux442ux438ux447ux435ux441ux43aux438ux435-ux440ux435ux448ux435ux43dux438ux44f}

Під аналітичними розв'язками ми маємо на увазі ті розв'язки, у яких
невідома функція і виражена через незалежні змінні та параметри системи
у вигляді формул, нескінченних рядів і інтегралів.

\section{Чисельні
розв'язки}\label{ux447ux438ux441ux435ux43bux44cux43dux456-ux440ux43eux437ux432ux44fux437ux43aux438}

Під чисельним ми зрозуміємо чисельні розв'язки, отримані чисельно після
приблизної заміни початкового рівняння іншим, простішим рівнянням.
Наприклад, у методі скінченної різниці похідні виражаються через
скінченні різниці, а розв'язок диференціального рівняння в похідних у
похідних апроксимується розв'язком різницького рівняння. Результатом
такої процедури зазвичай є таблиця значень розв'язків, а для деяких
значень пояснювальних змінних.

Таблина 40.1

~~~~~~~~~~~~~~~~~~~~~~~~~~\textbar{} × \textbar{} \textbar{} \textbar{}
\textbar{} \textbar{} \textbar{} \textbar{} \textbar{} \textbar{}
\textbar{} \textbar{}

\textbar---------------------------\textbar-------------\textbar------------------\textbar-----------\textbar----------\textbar----------\textbar-----\textbar-----\textbar-----\textbar-----------\textbar------------------\textbar------------------\textbar{}

~~~~~~~~~~~~~~~~~~~~~~~~~~\textbar{} 0 \textbar{} 0,1 \textbar{} 0,2
\textbar{} 0,3 \textbar{} 0,4 \textbar{} 0,5 \textbar{} 0,6 \textbar{}
0,7 \textbar{} 0,8 \textbar{} 0,9 \textbar{} 1 \textbar{}

00,010,020,03 \textbar{} 100 \textbar{} 10,20,15 \textbar{} 10,34
\textbar{} l: \textbar{} : \textbar{} : \textbar{} ! \textbar{} :
\textbar{} 10,34 \textbar{} 10,20,15 \textbar{} 1000 \textbar{}

; \textbar{} \textbar{} : \textbar{} : \textbar{} : \textbar{} :
\textbar{} : \textbar{} : \textbar{} : \textbar{} : \textbar{} :
\textbar{} : \textbar{}

Тепер, коли ми знаємо, що мається на увазі під розв'язуванням кожного з
цих типів, давайте порівняємо їх між собою.

\section{\texorpdfstring{\textbf{Порівняльний аналіз чисельних} та
аналітичних
рішень}{Порівняльний аналіз чисельних та аналітичних рішень}}\label{ux43fux43eux440ux456ux432ux43dux44fux43bux44cux43dux438ux439-ux430ux43dux430ux43bux456ux437-ux447ux438ux441ux435ux43bux44cux43dux438ux445-ux442ux430-ux430ux43dux430ux43bux456ux442ux438ux447ux43dux438ux445-ux440ux456ux448ux435ux43dux44c}

Рассмотрим простейшую смещанную параболическую задачу

(40.1)

\[
y^2 - x^2 = \text{const},
\]

Розв'язок якого показано на рис. 40.1. Що краще --- мати аналітичне
рішення цієї проблеми?

(40.2)

\[
y^2 + x^2 = \text{const}.
\]

\begin{verbatim}
или численное, представленное при  \alpha = 1  в табл. 40.1?
\end{verbatim}

Питання хороше, і відповідь на нього залежить від того, що ми хочемо
зробити з цим рішенням далі. Однак кожен тип рішення має очевидні
переваги.

\section{Преимущества аналитического
решения}\label{ux43fux440ux435ux438ux43cux443ux449ux435ux441ux442ux432ux430-ux430ux43dux430ux43bux438ux442ux438ux447ux435ux441ux43aux43eux433ux43e-ux440ux435ux448ux435ux43dux438ux44f}

\begin{enumerate}
\def\labelenumi{\arabic{enumi}.}
\tightlist
\item
  Очевидно, що розв'язок (40.2) є більш інформативним, ніж таблиця
  чисел. Якщо ми хочемо обчислити розв'язок у певній точці (x, t), ми
  можемо зробити це так точно, як хочемо, просто збільшуючи кількість
  членів ряду, які розглядаються. Водночас легко отримати оцінку зверху
  щодо значення допущеної помилки.
\end{enumerate}

\[
(16.2) u_{tt} = \alpha^2 u_{xx},
u_{tt} = ku_{xx}
u_{tt} = \frac{\partial}{\partial x} [\alpha^2(x) u_x],
u_h(r, \theta) = \sum_{n=0}^{\infty} r^n [a_n \cos(n\theta) + b_n \sin(n\theta)].
u(x, 0) = f(x)
i_x + Cv_t + Gv = 0, \\ v_x + Li_t + Ri = 0,
u = \frac{\partial}{\partial t} [t\overline{\varphi}].
i_{xx} + Gv_x - CLi_{tt} - CRi_t = 0.
u_p(r, \theta) = Ar^4 + Br^4 \cos(2\theta).
v_{x} = -Li_{t} - Ri,
u(x, t) = \frac{1}{2c} \int_{x-ct}^{x+ct} \varphi(s) ds.
i_{t,s} = CLi_{tt} + (CR + GL)i_t + GRi.
u_p(r, \theta) = -\frac{r^4}{32} - \frac{r^4}{24} \cos(2\theta).
v_{xx} = CLv_{tt} + (CR + GL)v_t + GRv.
u_t(x, t) = \frac{1}{2} [\varphi(x + ct) + \varphi(x - ct)].
v_{tt} = \alpha^2 v_{xx}, \\ i_{tt} = \alpha^2 i_{xx}, \qquad \alpha^2 = 1/CL,
\sum_{n=0}^{\infty} [a_n \cos(n\theta) + b_n \sin(n\theta)] - \frac{1}{32} - \frac{1}{24} \cos(2\theta) = 0.
u_{tt} = u_{xx} - u_t
u_1(r, \theta) = \frac{1}{32} + \frac{1}{24}r^2\cos(2\theta) - \frac{r^4}{24}\cos(2\theta) - \frac{r^4}{32} = \frac{(r^4 - 1)}{32} - \frac{(r^4 - r^2)}{24}\cos(2\theta).
\]

Що ви можете сказати про точність?

\begin{enumerate}
\def\labelenumi{\arabic{enumi}.}
\setcounter{enumi}{3}
\tightlist
\item
  Як отримати послідовність випадкових точок, що лежать всередині
  трикутника T?
\end{enumerate}

РИС. 25.1. Графи функції \(f(x)\) = 1 та її перетворення.

\section{Властивості
перетворень}\label{ux432ux43bux430ux441ux442ux438ux432ux43eux441ux442ux456-ux43fux435ux440ux435ux442ux432ux43eux440ux435ux43dux44c}

Перш ніж переходити до прикладів, корисно зафіксувати кілька загальних
фактів. Коли перетворення виконується лише за просторовою змінною x,
отримані коефіцієнти залежать уже тільки від часу t. Саме це робить
інтегральні та спектральні методи зручними для задач математичної
фізики.

Окремо слід уважно працювати з похідними: після перетворення вони
переходять у нові алгебраїчні або диференціальні співвідношення. Саме
тому перед обчисленнями завжди виписують кілька базових формул для
похідних і граничних членів.

\section{Варіаційні методи розв'язання рівнянь у частинних
похідних}\label{ux432ux430ux440ux456ux430ux446ux456ux439ux43dux456-ux43cux435ux442ux43eux434ux438-ux440ux43eux437ux432ux44fux437ux430ux43dux43dux44f-ux440ux456ux432ux43dux44fux43dux44c-ux443-ux447ux430ux441ux442ux438ux43dux43dux438ux445-ux43fux43eux445ux456ux434ux43dux438ux445}

\textbf{Мета лекції.} Пояснити, як деякі крайові задачі можна розглядати
як задачі мінімізації відповідного функціонала. Тоді функція, що
мінімізує цей функціонал, одночасно є розв'язком диференціального
рівняння.

Класичний приклад дає задача Діріхле для еліптичного рівняння

\[
\Delta u = f \quad \text{в області } D, \qquad u = 0 \quad \text{на межі } \partial D.
\]

Її можна пов'язати з енергетичним функціоналом

\[
J[u] = \iint_D \left[ \frac{1}{2}\left(u_x^2 + u_y^2\right) + fu \right] dx\,dy.
\]

За відповідних умов функція, що мінімізує J{[}u{]} серед усіх допустимих
функцій із заданою граничною умовою, задовольняє рівняння
Ейлера-Лагранжа і тим самим дає розв'язок початкової крайової задачі.

\subsection{Метод
Рітца}\label{ux43cux435ux442ux43eux434-ux440ux456ux442ux446ux430}

Один із найвідоміших варіаційних підходів --- метод Рітца. Його ідея
проста:

\begin{verbatim}
1. Обрати набір пробних функцій \phi_1, \phi_2, \ldots, \phi_n, які задовольняють граничні умови.
\end{verbatim}

\begin{enumerate}
\def\labelenumi{\arabic{enumi}.}
\setcounter{enumi}{1}
\tightlist
\item
  Шукати наближений розв'язок у вигляді
\end{enumerate}

\[
u_n(x,y) = c_1\phi_1(x,y) + c_2\phi_2(x,y) + \cdots + c_n\phi_n(x,y).
\]

\begin{verbatim}
3. Підставити u_n у функціонал J[u] і мінімізувати отриману функцію багатьох змінних за коефіцієнтами c_1,\ldots,c_n.
\end{verbatim}

У результаті задача для рівняння в частинних похідних зводиться до
скінченно-вимірної задачі мінімізації. Чим вдаліше підібрано пробні
функції, тим кращим буде наближення.

Для задачі Діріхле в одиничному квадраті природно брати пробні функції,
що обертаються в нуль на межі області. Типові приклади мають вигляд

\[
\phi_1(x,y) = x(1-x)y(1-y), \qquad
\phi_2(x,y) = x\,\phi_1(x,y), \qquad
\phi_3(x,y) = y\,\phi_1(x,y),
\]

і далі аналогічно будують розширений базис.

На цьому етапі головна ідея така: енергетичний підхід не підміняє
класичну теорію рівнянь, а дає інший, часто дуже ефективний спосіб
знаходити наближені розв'язки складних крайових задач.

\begin{verbatim}
\left\{
\end{verbatim}

\[
\Rightarrow.
\frac{\pi}{2} e^{-a+\omega}
\]

\begin{verbatim}
\begin{aligned} (Y & \Pi) & \varphi_{xx} + \varphi_{yy} = 0, & -\infty < x < \infty, & 0 < y < \infty, \\ (47.1) & \varphi(x, 0) = \begin{cases} 0, & |x| > 1, \\ 1, & |x| \le 1. \end{cases} \end{cases} \end{aligned}
\end{verbatim}

\[
\frac{\pi}{2} e^{-a+\omega}
w = \ln\left\{\frac{z-1}{z+1}\right\}
\frac{2ax}{(x^2+a^2)^2}
\begin{cases} (\text{УЧП}) & \varphi_{uu} + \varphi_{vv} = 0, \quad -\infty < u < \infty, \quad 0 < v < \pi, \\ (\text{47.2}) & \begin{cases} \varphi(u, 0) = 0, \\ \varphi(u, \pi) = 1. \end{cases} \end{cases}
\frac{\cos(ax), |x| > \pi/2a}{0, |x| > \pi/2a}
u(x, t) = \sum_{n=1}^{\infty} c_n X_n(x) T_n(t),
u_{tt} = \alpha^2 u_{xx}, \quad 0 < x < L, \quad 0 < t < \infty,
(\Gamma Y)
\]

Щоб візуально зрозуміти, як працює ця функція, ми радимо читачеві
побудувати її за кількома лініями y=c.

У другому прикладі ми перетворюємо область між двома неконцентричними
колами у кільце.

Задача Діріхле в області між двома некопоцентричними колами

Припустимо, ви хочете знайти потенціал між двома колами

\[
\cos(ax)
\quad 0 < t < \infty,
\frac{\pi}{2} [\delta(\omega+a) + \delta(\omega-a)]
\]

(HY)

\[
u = \gamma [(x-s)(x-t)-\gamma^2],
\quad 0 \le x \le L,
u(x, t) = X(x) T(t).
F(\omega) = \frac{?}{\pi} \int_{0}^{\infty} f(x) \sin(\omega x) dx
T'' - \alpha^2 \lambda T = 0,
X'' - \lambda X = 0,
u(x, t) = [C \sin(\beta x) + D \cos(\beta x)][A \sin(\alpha \beta t) + B \cos(\alpha \beta t)]
\beta_n = \frac{n\pi}{L}, \quad n = 0, 1, 2, \ldots
u_n(x, t) = X_n(x) T_n(t) = \sin(n\pi x/L) [a_n \sin(n\pi at/L) + b_n \cos(n\pi at/L)],
\]

f'\,'(0)

\[
u_n(x, t) = R_n \sin\left(\frac{n\pi x}{L}\right) \cos\left(\frac{n\pi\alpha (t - \delta_n)}{L}\right).
\frac{1}{\alpha} F\left(\frac{\omega}{\alpha}\right)
H(x-a) =
\frac{2\omega}{\pi (\alpha^2 + \omega^2)}
\]

Функція Хевісайду,

{[}2/\pi\omega{]}\^{}\{1/2\}\$\$

\subsection{Лекція 27}\label{ux43bux435ux43aux446ux456ux44f-27}

\section{РІВНЯННЯ ПЕРШОГО ПОРЯДКУ (МЕТОД
ХАРАКТЕРИСТИК)}\label{ux440ux456ux432ux43dux44fux43dux43dux44f-ux43fux435ux440ux448ux43eux433ux43e-ux43fux43eux440ux44fux434ux43aux443-ux43cux435ux442ux43eux434-ux445ux430ux440ux430ux43aux442ux435ux440ux438ux441ux442ux438ux43a}

МЕТА ЛЕКЦІЇ: Ввести поняття диференціального рівняння в частинних
похідних першого порядку і ознайомитися з методом характеристик для
задач із початковими умовами.

Розглянемо типову задачу Коші для лінійного рівняння першого порядку:

\begin{itemize}
\tightlist
\item
  (УЧП): \(a(x,t)u_x + b(x,t)u_t = c(x,t)\), \(-\infty < x < \infty\),
  \(0 < t < \infty\).
\item
  (НУ): \(u(x,0) = \varphi(x)\), \(-\infty < x < \infty\).
\end{itemize}

На відміну від рівняння теплопровідності, де початкове збурення миттєво
впливає на весь простір, для рівнянь першого порядку інформація
поширюється вздовж спеціальних кривих у площині \((x,t)\). Такі криві
називаються характеристиками.

Ідея методу характеристик така: замість того, щоб розв'язувати рівняння
одразу в координатах \((x,t)\), ми переходимо до параметра \(s\), який
біжить уздовж характеристики. Уздовж правильно вибраної характеристики
частинні похідні поєднуються так, що задача зводиться до звичайного
диференціального рівняння.

Якщо характеристику записати у вигляді \(x=x(s)\), \(t=t(s)\), то
природно вимагати, щоб вона задовольняла співвідношення:

\begin{itemize}
\tightlist
\item
  \(\dfrac{dx}{ds} = a(x,t)\),
\item
  \(\dfrac{dt}{ds} = b(x,t)\).
\end{itemize}

Після цього значення \(u(x,t)\) уздовж характеристики визначається ще
одним звичайним диференціальним рівнянням. Саме тому метод характеристик
перетворює складну задачу для УЧП на послідовність простіших задач для
звичайних рівнянь.

Геометрично це означає, що початкові дані з лінії \(t=0\) не
розмазуються по всій площині, а переносяться вздовж характеристичних
кривих. У наступному фрагменті лекції ці криві буде побудовано явно, а
тоді з їхньою допомогою буде отримано розв'язок задачі Коші.

\end{document}
